\documentclass[11pt]{article}
\usepackage[T1]{fontenc}
\usepackage{lmodern}
\usepackage[margin=0.88in,headheight=14pt]{geometry}
\usepackage{amsmath,amssymb,amsthm,mathtools}
\usepackage{booktabs,longtable,array}
\usepackage{microtype,xcolor,enumitem,fancyhdr}
\usepackage{hyperref,xurl}
\hypersetup{colorlinks=true,linkcolor=blue!45!black,urlcolor=blue!45!black,
 pdftitle={Two Routes to Smaller Diophantine Consistency Equations},
 pdfauthor={AI-assisted research prepared for Vladimir Reshetnikov}}
\newtheorem{theorem}{Theorem}[section]
\newtheorem{lemma}[theorem]{Lemma}
\newtheorem{proposition}[theorem]{Proposition}
\theoremstyle{remark}
\newtheorem{remark}[theorem]{Remark}
\newcommand{\Z}{\mathbb Z}
\newcommand{\N}{\mathbb N}
\newcommand{\Zpos}{\mathbb Z_{>0}}
\newcommand{\Con}{\operatorname{Con}}
\newcommand{\ZFC}{\mathrm{ZFC}}
\newcommand{\rad}{\operatorname{rad}}
\newcommand{\vp}{\nu_p}
\newcommand{\yb}{\overline Y}
\newcommand{\PellX}{\mathsf X}
\newcommand{\PellY}{\mathsf Y}
\pagestyle{fancy}
\fancyhf{}
\lhead{\small Diophantine consistency}
\rhead{\small Weighted transport and Pell witnesses}
\cfoot{\thepage}
\setlength{\parskip}{4pt}
\setlength{\parindent}{0pt}
\setlist{itemsep=2pt,topsep=4pt}
\allowdisplaybreaks

% BEGIN retained input: direct_metrics.tex
\newcommand{\DirectVariables}{83}
\newcommand{\DirectDegree}{624}

% END retained input: direct_metrics.tex


\title{\textbf{Two Routes to Smaller Diophantine\\Consistency Equations}\\[6pt]
 \large Weighted transport and reuse of Pell witnesses}
\author{Research prepared for Vladimir Reshetnikov and ProveIt\\
 \small AI-assisted mathematical development with separate internal audits}
\date{25 October 2026}
\begin{document}
\maketitle

\begin{abstract}
We construct an explicit equation in \(\DirectVariables\) unrestricted
integers of exact degree \(\DirectDegree\), through a direct numerical
proof calculus, and an explicit equation in ten strictly positive integers
of exact degree \(190\). Using the named inherited logical calibrations,
ZFC proves that each equation's insolubility in its stated domain is
equivalent to \(\Con(\ZFC)\).
The direct construction reuses the quotient of a large Pell solution as
its modular center witness, then applies certified acyclic unit
elimination. It uses neither K215 simulation nor Jones's universal
polynomial. Its factored evaluator uses \(1{,}789\) binary operations
and integer literals of absolute value at most \(64\).
The ten-variable construction changes a positive transport
weight so that a center for an entire Pell product works in a newer
free-quotient chart. It retains the native/K215 coefficient ancestry and
uses \(118\) live binary arithmetic operations. A companion has degree
\(202\) and \(99\) operations.

The archive supplies literal equations, all fixed coefficients, finite
reconstruction, exact algebra checks, and reproducible code.
The proofs are ordinary mathematical arguments composed with identified
inherited interfaces; no new end-to-end Lean/Rocq formalization is claimed.
The results establish neither consistency nor an actual zero, and do
not establish a global simplicity minimum.
\end{abstract}
\clearpage
\tableofcontents
\newpage

\section{Results, comparison, and logical scope}

Fix the effective syntactic presentation of ZFC used by the repository's
refutation calibration. Its consistency sentence says that no finite proof
in that presentation derives a contradiction. Every quantified integer
below is an integer internal to the ambient theory; no restriction to
externally standard proof codes is made.

There are two main outputs, with different domains and costs.
The direct route produces the polynomial
\(D_{\DirectVariables}\) described in Section~\ref{sec:direct-results} and
fully specified in Appendix~\ref{app:direct83-definition} and its literal
electronic appendix. It has
\(\DirectVariables\) unrestricted integer inputs and exact degree
\(\DirectDegree\). The second route gives the following theorem.

\begin{theorem}[Ten positive variables]\label{thm:principal}
Let \(P_{190}\) be the polynomial defined in Section~\ref{sec:equation},
with the four fixed integers specified in Appendix~\ref{app:constants}.
Then
\[
 \ZFC\vdash
 \left(\neg\exists\boldsymbol x\in\Zpos^{10}\;
                      P_{190}(\boldsymbol x)=0\right)
 \ \longleftrightarrow\ \Con(\ZFC).
\]
Its exact total degree is \(190\). Its supplied literal circuit uses
\(73\) multiplications, \(24\) additions, and \(21\) subtractions.
\end{theorem}

The proof composes the new arithmetic arguments with the inherited
calibration identified in Section~\ref{sec:boundary}.
The coefficient specification is independent of the truth of consistency.
For the direct route the analogous theorem quantifies over
\(\Z^{\DirectVariables}\); that domain distinction is substantive.

\begin{center}
\begin{tabular}{lrrr}
\toprule
Ten-positive-variable construction & Degree & Operations & Root scale\\
\midrule
Source 172, free quotient and inverse correction & 214 & 87 & 1\\
Source 178, whole-product center in older chart & 208 & 116 & 4\\
Incoming Affine Horizons, centered free chart & 206 & 90 & 1\\
This report, quotient-center companion & 202 & 99 & 2\\
\textbf{This report, weighted whole-product center} & \textbf{190} &
 \textbf{118} & \textbf{4}\\
\bottomrule
\end{tabular}
\end{center}

The comparison snapshot is the immutable repository commit
\path{cdfea7d90176775d74c1c30afbe49fa78ade9897}.
The incoming archives at that snapshot were inspected as well as the
placed packages. Thus degree \(206\), rather than the placed
degree \(208\), is the relevant preceding ten-variable degree benchmark
\cite{source172,source178,incomingAffine}.
The drop \(206\to190\) is \(16\) degrees, approximately \(7.77\%\).
The operation count rises \(90\to118\); neither point dominates the
other under both measures.

The word \emph{simplest} requires a comparison rule. The same repository
already contains nine-positive-variable equations of much higher degree,
ten-unrestricted-integer equations using a different reduction, and
quartics with more variables and large fixed coefficients.
For example, package 101 supplies a nine-positive degree-\(997{,}810\)
equation, and package 82 a quartic in \(37\) unrestricted integers.
The new results therefore establish measured improvements in identified
families, not minimum variable counts or global optimality.
For the ten-positive route the coefficient cost is explicit:
the largest port has \(592{,}884\) bits. Counting multiplication by that
integer as one operation is not a bit-complexity assertion.

To obtain a ten-variable equation over nonnegative integers without
changing total degree, replace every input \(x_j\) by \(x_j+1\).
That translated circuit has ten additional additions. The positive-domain
theorem cannot simply be relabeled as an unrestricted-integer theorem.

\section{The ten-positive-variable polynomial}\label{sec:equation}

The independent inputs and their entire domain are
\[
 (J,F,\alpha,i,T,w,s,n,\eta,\zeta)\in\Zpos^{10}.
\]
The symbols \(\beta,K_0,M_C,M_F\) denote the four fixed integers of
Appendix~\ref{app:constants}. Every other letter in this section is a
polynomial abbreviation, not an additional unknown.

\subsection{Common arithmetic registers}
Define
\begin{align}
 h&=\beta J,&q&=h+1,&C&=q-F-\alpha, \nonumber\\
 X&=qw,&Y&=s,&\yb&=qY, \nonumber\\
 k&=\eta+\zeta,&c&=kY+\eta,&
 A&=\yb(X+1)+1, \label{eq:front}\\
 \Delta&=A^2-1,&H&=4A-5,&v&=X\yb^2,\nonumber\\
 \ell&=vk,&E&=1+\ell(\ell+k).&&\nonumber
\end{align}
The remaining shared registers are
\begin{align}
 R&=[h(q-F)+\alpha](q^2-1)+(M_C+qM_F)J, \nonumber\\
 B&=(K_0+w)C-F,&U&=ic^2, \nonumber\\
 a&=\Delta U^2E,&F_0&=1+a,&G&=EF_0, \label{eq:pellregs}\\
 f&=1+2a,&Q&=4\Delta aF_0,&b&=cT-R,\nonumber\\
 V&=fb-c,&K&=X(3c-2X)+2,&&\nonumber\\
 e_6&=6(k-R-1)-\ell(k^2-4),&
 D&=6q^7,&M&=hDH.\nonumber
\end{align}
The subscript on \(K_0\) distinguishes the fixed coefficient from the
polynomial \(K\). The numeral \(6\) in \(D\) is part of the literal equation.

\subsection{The center and the transport}
Set
\begin{align}
 e_2&=(k\yb)^2,\nonumber\\
 S&=(k\yb)\bigl[\ell+k-(k\yb)(\yb+1)\bigr],\nonumber\\
 \theta&=U^2b,\qquad
 R_c=\theta(3E+e_2)-(b-c),\nonumber\\
 S_c&=4E^2+3S^2+e_2^2,\label{eq:center}\\
 \Omega&=U\bigl[2SR_c+A(4EV-\theta S_c)\bigr],\nonumber\\
 L&=AU^3b^2.\nonumber
\end{align}
The useful folded evaluation is
\[
 \Omega_{\mathrm f}
   =U\bigl[2SR_c+A\{4EV-\theta(S_c-bB)\}\bigr].
 \tag{F}
\]
It satisfies the exact identity
\[
 \Omega_{\mathrm f}=\Omega+LB.
 \tag{F1}
\]
One can therefore compute the carrier directly as
\begin{align}
 Z&=hD\bigl[H(n-\Omega_{\mathrm f})+K\bigr]
      -H\bigl[DB+h(e_6-6q^4s)\bigr],             \label{eq:Z}\\
 P_{190}&=Q\bigl(Z^2-16M^2GV^2\bigr)+16M^2G-Z^2.
                                                        \label{eq:P}
\end{align}
\begin{center}
\fbox{\(\displaystyle
 P_{190}(J,F,\alpha,i,T,w,s,n,\eta,\zeta)=0,
 \quad J,F,\alpha,i,T,w,s,n,\eta,\zeta>0. \)}
\end{center}
Equations \eqref{eq:front}--\eqref{eq:P}, together with the coefficient
appendix, specify one ordinary polynomial with integer coefficients.
Quotients, modular arithmetic, Pell sequences, and roots used below belong
only to the proof. They are not primitives in this polynomial.

\section{What is inherited and what changes}\label{sec:boundary}

Source 172 supplies the free chart \(Y=s,\ \yb=qY\), the inverse-Lucas
correction \(e_6\), the denominator \(6q^7\), and the signed analysis that
permits auxiliary Pell index three. Its complete local proof and the
selected underlying proof files are retained in the archive.
Source 178 supplies the center for the whole product and its
polynomial-quotient identity. The incoming Affine Horizons archive supplies
the degree-\(206\) comparison \cite{source172,source178,incomingAffine}.

The new contribution is the positive weight \(L=AU^3b^2\), its independent
bound in the free chart, the folding identity (F1), and the literal
degree-\(190\) equation. These ingredients are needed together.
An older unweighted \(b^2\) estimate relied on excluding an auxiliary index
that the free chart allows. Importing that exclusion would invalidate
the argument.

After root extraction and sign recovery, the inherited Pell authentication
and ratio projector recover
\[
 X=2^R,\qquad q^3\mid Y,\qquad C>0,\qquad
 \frac{(K_0+w)C-F}{h}\in\N.
 \tag{N}
\]
The entire calibrated four-number tuple then reconstructs the native
words at the source base \(Y(X+1)+2\); the inherited native/K215 refutation
calibration produces a finite syntactic ZFC refutation.
Conversely, a refutation gives canonical positive outer data, compatible
Pell witnesses, and the completeness bound used below.
Neither direction permits replacing the actual four integers by arbitrary
integers satisfying the same size inequalities.

The direct construction in Section~\ref{sec:recycle} uses a different
calibration: the numerical C222 calculus, then ZF refutation and ZFC
refutation through the named
\texttt{KEndpoint.c222\_iff\_zf\_refutation} and
\texttt{ConstructibleChoice.zf\_refutation\_iff\_zfc\_refutation}
interfaces. It does not use the four-number tuple.

These are mathematical dependency boundaries. No newly compiled complete
Lean/Rocq import closure or object-language ZFC derivation is claimed.
The new arguments use integer algebra, finite factorization, induction,
and the inherited theorems, and can be carried out internally in ZFC.
Finite checks validate emitted arithmetic; they do not prove unbounded
semantic assertions by sampling. The archive's \texttt{calibration/} folder
also records an inspection of the actual unconditional endpoint declarations,
their mathematical source dependencies, and selected archived assumption logs.
It retains precise source pins and distinguishes those historical logs from
a fresh proof-assistant build.

\section{Soundness of the degree-190 equation}

\subsection{Root extraction before sign recovery}
At every positive input, \(q\ge32,\ X\ge q,\ Y\ge1,\ k\ge2,\ c\ge3\),
\(U\ge c^2\), and
\[
 A\ge32(32+1)+1=1057.                              \tag{B0}
\]
The identities
\[
 f^2=1+4aF_0,\qquad
 Q=\Delta(f^2-1)=4\Delta^2U^2G,\qquad
 \gcd(E,F_0)=1                                    \tag{I}
\]
are independent of the carrier. Since \(f>2c\), \(V=fb-c\)
is nonzero for every integral \(b\).
Suppose \(P_{190}=0\). Rewriting the final norm gives
\[
 (Q-1)Z^2=16M^2G(QV^2-1).                         \tag{R}
\]
A prime dividing \(G\) divides neither \(Q-1\) nor \(QV^2-1\).
Valuations in (R) make \(G\) a square. Coprimality in (I)
gives positive integers \(\tau,f_0,g\) with
\[
 E=\tau^2,\quad F_0=f_0^2,\quad G=g^2,\quad g=\tau f_0. \tag{R1}
\]
The factor \(16\) changes no valuation parity.
The inherited adjacent-coefficient Pell descent applies initially to the
rational number \(|Z|/(4Mg)\) and proves it integral. Thus
\[
 QV^2-(Q-1)y^2=1,\qquad Z/M=\epsilon\,4gy,\qquad
 y\in\Zpos,\quad\epsilon\in\{-1,1\}.              \tag{R2}
\]
One does not assume numerator divisibility in advance.

The inherited signed all-index folding, applied to
\[
 I=2Ug,\quad f^2-\Delta I^2=1,\quad
 Q=(\Delta I)^2,\quad c^2\mid I,\quad V\equiv-c\pmod f,
\]
gives \(c=\PellY_p(A)\) for positive odd \(p\).
As \(c\ge3\), \(p\ge3\), and
\[
 c\ge4A^2-1=4\Delta+3,\qquad c>X,\qquad K>0.      \tag{R3}
\]
No sign of \(B,C,R,b,V\), and no divisibility of \(Y\) by \(q\), has
been used.
The same parent analysis proves that \(b,V\) are nonzero with the same
sign, including the auxiliary-index-three case, and
\[
 |b|>4\Delta f-1>3\Delta f,\qquad
 y<\tfrac32|V|,\qquad |V|<\tfrac32f|b|,\qquad f>6g. \tag{R4}
\]
These signed bounds are the appropriate input to the new argument.

\subsection{An absolute bound for the polynomial center}
For proof only, replace \(A\) by a formal indeterminate \(t\), keeping
\(\yb,k,U,b,c\) fixed. Put
\begin{align*}
 E(t)&=e_2t^2+e_1t+e_0,\\
 e_1&=k^2\yb(1-2\yb-2\yb^2),\\
 e_0&=1+k^2\yb(\yb+1)\bigl[\yb(\yb+1)-1\bigr],\\
 V(t)&=\bigl[1+2(t^2-1)U^2E(t)\bigr]b-c,\qquad
 F(t)=E(t)V(t).
\end{align*}
Let \(\mathcal Q_j\) be the ordinary polynomial quotient of \(F(t)\) by
\(t^j\), evaluated at \(A\). Exact algebra gives
\[
 \Omega=U\left(4AEV-2\mathcal Q_1-\frac{\mathcal Q_3}{2}\right).
 \tag{C1}
\]
The quotient \(\mathcal Q_3\) is coefficientwise even: the term
\(E(t)(b-c)\) has degree at most two and the remaining term has factor
\(2U^2b\). Equation (C1) is proof notation for the division-free center
in \eqref{eq:center}, not numerical polynomial division in the equation.

Let \(F^\#(A)\) mean evaluation at positive \(A\) after replacing every
coefficient by its absolute value. The bound
\(\delta=(\yb+1)/A<1/16\) follows from \(\yb,X\ge32\).
Using \(|e_1|<2e_2(\yb+1)\) and \(0<e_0\le e_2(\yb+1)^2\) gives
\[
 E^\#(A)<2E,\qquad
 F^\#(A)<2E(|b|+c)+8(A^2+1)U^2E^2|b|.
\]
Since \(c+1<2U^2E\) and \(|V|>\Delta U^2E|b|\),
\[
 \frac{F^\#(A)}{|EV|}<8+\frac{20}{\Delta}<10.      \tag{C2}
\]
Deletion of low coefficients yields
\(|\mathcal Q_j|\le F^\#(A)/A^j\). Hence
\[
 |\Omega-4AUEV|
 <10U|EV|\left(\frac2A+\frac1{2A^3}\right)<AUE|V|.
\]
We obtain
\[
 \operatorname{sgn}(\Omega)=\operatorname{sgn}(V),\qquad
 3AUE|V|<|\Omega|<5AUE|V|.                        \tag{C3}
\]
No premature assumption \(V>0\) is hidden in this estimate.

\subsection{The weighted transport inequality}
\begin{lemma}\label{lem:transport}
At every positive zero of \(P_{190}\),
\[
 E<U^2,\qquad 4gy+|\Omega|<\tfrac12L,\qquad L=AU^3b^2>0.
\]
\end{lemma}
\begin{proof}
From \(v=X\yb^2\), \(A=\yb(X+1)+1\), and (R3),
\[
 v+1<A^2<c,\qquad k<c.
\]
These quantities are integers, so
\[
 E=1+k^2v(v+1)\le1+(c-1)^4<c^4\le U^2.
\]
Also \((AUE)^2-G=U^2E^2-E>0\), so \(g<AUE\).
Combine this with (C3) and (R4):
\begin{align*}
 4gy+|\Omega|
  &<6g|V|+5AUE|V|\\
  &<11AUE|V|<\frac{33}{2}AUEf|b|\\
  &<\frac{11E}{2\Delta U^2}\,AU^3b^2
   <\frac{11}{2\Delta}L<L/2.
\end{align*}
The last inequality uses \(\Delta>11\).
All estimates precede recovery of the sign of \(B\).
\end{proof}

The weight fits both the signed estimate and a factor already present in
the center. Identity (F1) absorbs it into the literal evaluator.

\subsection{Recover the arithmetic sign}
The following are rational abbreviations used only in the proof:
\[
 e^*=e_6/6,\quad \rho=e^*/q^7,\quad
 \sigma=s/q^3,\quad t=B/h,\quad\kappa=K/H,\quad
 z=Z/M,\quad W_c=1+hL.
\]
Expansion gives the exact identity
\[
 z=n+\kappa+\sigma-W_ct-\rho-\Omega.              \tag{T}
\]
We have \(\kappa,\sigma>0\). Since \(k\ge2\), the inverse correction obeys
\[
 e^*\le k-R-1<b,\qquad \rho<|b|/q^7.             \tag{T1}
\]
This remains valid for negative \(R\).
If \(B\le-1\), then \(W_ct<-L\). As \(L\ge|b|\),
Lemma~\ref{lem:transport} and (R2) give
\[
 n=\epsilon4gy+\Omega-\kappa-\sigma+W_ct+\rho
       <L/2-L+|b|/q^7<0.
\]
This contradicts the independent condition \(n>0\).
Thus \(B\ge0\), and then \(C>0\).

\subsection{Authentication and native recovery}
Only now may one invoke the calibrated mask estimate \(3q<R<q^4-2\).
It gives \(b,V>0\); the unchanged signed index argument yields
\(p=R\), \(R\) odd, \(q\) even, and \(J\) odd.
The carrier changes do not affect
\(V\equiv-c\pmod f\) or \(V\equiv-R\pmod c\).

The center shift \(M\Omega\) is an integer multiple of \(D\) and \(H\),
and \(L\) is an integer polynomial. The parent numerator reductions
and fractional-part argument therefore remain valid, in this order:
\begin{enumerate}
\item The first Pell pair gives \(k=2\PellY_m(1+2v)\).
Literal cancellation of \(6\) supplies \(q^4\mid e^*\).
\item The retained cubic-difference identity, including its prime-\(3\)
valuation argument, yields \(m=(R+1)/2\).
One does not yet infer \(q^7\mid e^*\).
\item Reduction modulo \(H\) gives \(H\mid6hK\).
The \(\lambda=q\) ratio theorem and universal projector give \(X=2^R\).
Since \(X=qw\), \(q\) is a power of two.
\item The stronger resulting valuation of \(v=Xq^2Y^2\) yields
\(q^7\mid e^*\). Now \(\rho,\kappa\) are integral, and (T) gives
\(q^3\mid Y\) and \(h\mid B\).
\end{enumerate}
These are the calibrated endpoint (N). Fresh native reconstruction
produces a finite ZFC refutation. This completes soundness.

\section{Completeness and the consistency equivalence}
Start with a finite syntactic ZFC refutation. The inherited construction
gives canonical positive outer data, compatible positive Pell witnesses,
and a positive trace-linear coordinate \(T\), with
\[
 b,V,y>0,\quad t=B/h\in\N,\quad
 \rho,\kappa,\sigma\in\Z,\quad |\rho|+\kappa+\sigma<g. \tag{E}
\]
The last bound is proved independently of the carrier weight.
Retain every input except \(n\) and set
\[
 n=4gy+\Omega+(1+hL)t+\rho-\kappa-\sigma.          \tag{E1}
\]
It is an integer. Since \(\Omega>0\), \(t\ge0\), and \(y\ge1\),
(E) proves it positive. Equation (T) says \(Z/M=4gy\), and the final
norm vanishes.

A positive zero therefore exists if and only if a finite syntactic ZFC
refutation exists. Negation proves Theorem~\ref{thm:principal}.
The argument does not assume consistency or assert insolubility.

\begin{remark}
Completeness constructs canonical witnesses. It is not a pointwise
bijection of all old and new positive zero sets. The inverse of a polynomial
translation of a positive coordinate need not remain positive.
The separate signed soundness argument is essential.
\end{remark}

\section{Exact degree and a cheaper companion}

\subsection{The degree-190 certificate}
Let \(Z_0\) be the carrier with transport \(L\), before subtraction of
\(M\Omega\). Thus \(Z=Z_0-M\Omega\).
Regard \((A,\yb,k,U,b,c)\) as free variables of weights
\((4,2,1,5,4,2)\). Exact collection gives \(594\) nonzero terms in the
unfolded center defect, weighted degree \(94\), and top part
\[
 [\Omega^2-16GV^2]_{94}
       =8A^6\yb^8k^8U^6b^2(A^2-6\yb^4).         \tag{D1}
\]
Using the folded center instead gives \(679\) terms, with the same top
part; its additional transport terms have lower weight.
In the ten actual inputs,
\[
 \deg(Q,M,G,V,\Omega,Z_0)=(72,12,46,36,59,41).
\]
The cancellation-sensitive expansion is
\[
 P_{190}
 =Q\{Z_0^2-2MZ_0\Omega+M^2(\Omega^2-16GV^2)\}
       +16M^2G-(Z_0-M\Omega)^2.                  \tag{D2}
\]
The five degree bounds are \(154,184,190,70,142\); only the third term
can contribute in degree \(190\).
Put \(q_0=\beta J\), \(k_0=\eta+\zeta\). Substitution in (D1) gives
the entire top homogeneous form
\[
 [P_{190}]_{190}
  =9\,2^{11}q_0^{74}w^{18}s^{52}i^{10}k_0^{32}
       (q_0-F)^2(w^2-6s^2).                      \tag{D3}
\]
It is nonzero for the actual \(\beta>0\), proving exact degree.
The literal graph specifies the whole polynomial without expansion.

\subsection{A 99-operation companion}
Keep the common definitions \eqref{eq:front}--\eqref{eq:pellregs} and \(S\).
For this alternative only, put
\[
 L=bV,\quad C_2=2AUE-US,\quad
 Z=hD[H\{n-V(bB+C_2)\}+K]-H[DB+h(e_6-6q^4s)]
\]
and
\[
 P_{202}=Q(Z^2-4M^2GV^2)+4M^2G-Z^2.              \tag{C202}
\]
These definitions replace the principal carrier; there remain ten positive
inputs. The positive remainder
\[
 R_E=1+k^2\yb(\yb+1)\{\yb(\yb+1)-1\}
\]
satisfies \(E=AS+R_E\), \(0<S<E/A\), and
\(0<C_2<2AUE<3g\). The inherited bound gives \(|b|>16g\), so
\[
 2gy+C_2|V|<6g|V|<L/2.
\]
The same bootstrap and completeness construction work with root scale two.
The identity
\[
 C_2^2-4G=4U^2ER_E+U^2S^2-4E
\]
has degree \(34\). The output has exact degree \(202\), with top form
\[
 [P_{202}]_{202}
  =9\,2^{12}q_0^{80}w^{22}s^{56}i^{10}k_0^{32}(q_0-F)^2.
\]
Its evaluator uses \(60\) multiplications, \(21\) additions, and
\(18\) subtractions.

\section{Reusing a Pell quotient as a center witness}\label{sec:recycle}

The preceding construction retains a native coefficient calibration.
The following mechanism operates in the direct C222 proof-calculus branch
and removes a private witness by changing how a large center is chosen.
Use standard Pell polynomials
\[
 (a+\sqrt{a^2-1})^j
   =\PellX_j(a)+\PellY_j(a)\sqrt{a^2-1},
 \qquad a\ge2,\quad j\ge0.
\]
Their recurrences and addition laws are polynomial identities.

\begin{theorem}[The quotient supplies the center]\label{thm:reuse}
Let \(q,\mathcal M\) be positive odd integers with
\(\rad(\mathcal M)\mid q\), and let \(W_0\ge2\).
Set
\[
 W=q^2W_0,\quad \mathcal K=W^3(W+2),\quad
 A_0\equiv1+\frac{q^2}{2(q+1)}\pmod{\mathcal M}.
\]
There are arbitrarily large positive odd \(j\) for which
\[
 L_0=\PellX_{Wj}(W+1),\quad
 r=\frac{\PellY_{Wj}(W+1)}{W},\quad A_*=L_0^2+2r^2
\]
satisfy
\[
 r\in\Zpos\text{ is odd},\quad
 L_0^2-\mathcal K r^2=1,\quad A_*\equiv A_0\pmod{\mathcal M}.
\]
In particular \(A_*\ge L_0^2\), and these centers are cofinal.
\end{theorem}
The quotient \(r=2z+1\) is already a source witness.
No extra coordinate stores \(j\), a modular root, or a CRT choice.

\subsection{A local unit and a permutation of indices}
\begin{lemma}\label{lem:unit}
If \(p\) is odd and \(p^2\mid W\), then
\[
 b_0=\PellY_W(W+1)/W\in\Z,\quad b_0\equiv1\pmod p,\quad
 c_0=\PellX_W(W+1)\equiv1\pmod{p^6}.
\]
\end{lemma}
\begin{proof}
The recurrence gives \(\PellY_j(W+1)\equiv j\pmod W\).
With \(a=W+1\), \(d=W(W+2)\), binomial expansion gives
\[
 b_0=\sum_{2t+1\le W}
       \frac{\binom W{2t+1}}{W}\,a^{W-2t-1}d^t.
\]
The \(t=0\) term is \(1\) modulo \(p\). For \(t\ge1\),
\[
 \vp\left(\binom W{2t+1}/W\right)\ge-\vp(2t+1).
\]
Each whole term has valuation at least \(2t-\vp(2t+1)\ge1\).
Although the displayed coefficients can be rational, each whole term
is \(p\)-integral, and these valuation bounds are legitimate.
Thus \(b_0\equiv1\pmod p\). Also \(c_0\equiv1\pmod p\).
The norm identity \(c_0^2-1=dW^2b_0^2\) has right side divisible by
\(p^6\), while \(c_0+1\) is a unit. This proves the claim.
\end{proof}

\begin{lemma}[Index permutation]\label{lem:permutation}
If \(p\) is odd and \(c_0\equiv1\pmod{p^2}\), then, for \(a,j\ge0\),
\[
 \PellY_{j+p^a}(c_0)-\PellY_j(c_0)
       \equiv p^a\pmod{p^{a+1}}.                 \tag{Perm}
\]
Hence \(j\mapsto\PellY_j(c_0)\) is a permutation modulo every power of \(p\).
\end{lemma}
\begin{proof}
Put \(h=p^a\), \(d'=c_0^2-1\).
The elementary bound \(\vp\binom h t\ge a-\vp(t)\) and
\(\vp(d')\ge2\) give
\[
 \PellX_h(c_0)\equiv1\pmod{p^{a+1}},\qquad
 \PellY_h(c_0)\equiv p^a\pmod{p^{a+1}}.
\]
Every nonleading ordinate term has valuation at least
\(a+2t-\vp(2t+1)\ge a+1\).
The scalar expansion is similar; its leading power \(c_0^h\) is
\(1\) modulo \(p^{a+1}\).
The addition law
\[
 \PellY_{j+h}(c_0)=\PellX_h(c_0)\PellY_j(c_0)
                         +\PellX_j(c_0)\PellY_h(c_0)
\]
proves (Perm), using \(\PellX_j(c_0)\equiv1\pmod p\).
Modulo \(p\) the map is the identity. At the next lift, adding \(tp^a\)
changes the new base-\(p\) digit by \(t\). Induction proves bijectivity.
\end{proof}

\subsection{Proof of the center theorem}
If \(\mathcal M=1\), every odd positive \(j\) suffices.
Otherwise \(p^2\mid W\) for each \(p\mid\mathcal M\).
The composition law
\[
 r(j)=\PellY_{Wj}(W+1)/W=b_0\PellY_j(c_0)
\]
and the lemmas make \(j\mapsto r(j)\) a permutation modulo every prime
power dividing \(\mathcal M\). CRT gives a permutation modulo \(\mathcal M\).
For each such prime,
\[
 2(q+1)(\mathcal K+2)\equiv4\pmod p.
\]
The simple roots \(2,-2\) lift to all required prime powers by
\(d\mapsto d+tp^a\), since \(2d\) is invertible.
CRT supplies a unit \(d\pmod{\mathcal M}\) with
\[
 d^2\equiv2(q+1)(\mathcal K+2)\pmod{\mathcal M}.
\]
Choose \(j\) with \(r(j)\equiv q/d\pmod{\mathcal M}\).
Because \(\mathcal M\) is odd, this class has odd positive representatives;
adding \(2\mathcal M\) gives arbitrarily large ones.
The norm identity now gives
\[
 A_*=1+(\mathcal K+2)r^2
       \equiv1+\frac{q^2}{2(q+1)}\pmod{\mathcal M}.
\]
For odd \(W\), \(\PellY_W(W+1)/W\) is odd by reduction modulo two.
For even \(W\), all abscissae at \(W+1\) are odd, and
\(\PellY_{2t}=2\PellX_t\PellY_t\) proves
\(\nu_2(\PellY_W)=\nu_2(W)\). The quotient is again odd.
Finally \(\PellY_j(c_0)\) is odd for odd \(j\), regardless of \(c_0\)'s
parity. Positivity and cofinal growth follow from the positive recurrence.

\subsection{Why this identifies a new pair of witness roles}
The incoming cofinal-center archive already proves that an even cofinal
frame permits a separate witness \(A_*=L_0^2+2u^2\)
\cite{incomingCofinal}. Here the existing \(r\) replaces \(u\).
The changed frame lets that quotient cover the required residue classes.
The hypothesis matters: at \(W=3,p=3\),
\[
 \PellY_3(4)/3=21
\]
is not a local unit. Requiring merely \(p\mid W\) would not justify the proof.

In package 166 take \(W_0=2+\sum a_i^2\), using exactly its old affine atoms.
Replace
\[
 W-W_0=0,\qquad A_*-L_0^2-u^2-v^2=0
\]
by
\[
 W-q^2W_0=0,\qquad A_*-L_0^2-2(2z+1)^2=0.       \tag{Reuse}
\]
Delete the now private \(u,v\); all other source rows are unchanged.
Soundness first uses the affine radix row \(q=6t+3\), hence
\(W=q^2W_0\ge W_0\), and \(A_*\ge L_0^2\).
It does not assume \(\mathcal M\ne0\) before power recovery.
Completeness first constructs the semantic positive odd \(q,\mathcal M\)
with the required prime support, and then applies Theorem~\ref{thm:reuse}.

This is an equivalence of zero existence, not a map from every preselected
old center. The raw frame row has degree four, giving a degree-eight
sum of squares. A signed-frame refinement and exact elimination give
the stronger direct endpoints below.


% BEGIN retained input: direct_results.tex
\section{A direct equation in 83 unrestricted integers}\label{sec:direct-results}

The modular witness-reuse theorem leads to a concrete second construction.
Its source is the unmodified direct numerical-calculus polynomial in
package 166, followed by the two semantic changes just proved and explicit
integral coordinate elimination. The main endpoint is
\[
 D_{83}(x_0,\ldots,x_{82})
       =\sum_{j=0}^{30}R_j(x_0,\ldots,x_{82})^2.
 \tag{Direct}
\]
Appendix~\ref{app:direct83-definition} prints every free assignment, every
ordered substitution, and every residual. The electronic appendix also
supplies the corresponding literal arithmetic circuit and source coefficients.
All \(83\) inputs range independently over \(\Z\).
There are no hidden inequalities or extra witnesses.

\begin{theorem}[Direct C222 route]\label{thm:direct}
For the fixed polynomial (Direct), with the inherited C222/ZF/ZFC
refutation calibration identified in Section~\ref{sec:boundary},
\[
 \ZFC\vdash
 \left(\neg\exists\boldsymbol x\in\Z^{83}\;D_{83}(\boldsymbol x)=0\right)
       \longleftrightarrow\Con(\ZFC).
\]
The exact total degree is \(624\), and the preferred factored evaluator has
\(1{,}789\) live binary arithmetic operations. Its integer literal scalars
have absolute value at most \(64\).
\end{theorem}

This route does not simulate K215 or specialize Jones's universal
polynomial. It retains a specific direct proof calculus and its previously
identified consistency calibration. Compared with the \(2{,}428\)-operation \(92\)-integer, degree-\(692\) equation
\cite{source175}, the new point improves dimension \(92\to83\), degree
\(692\to624\), and operation count \(2{,}428\to1{,}789\) together.
The latter count uses the exact factorizations described below; the
canonical evaluator retained for the projection audit has \(2{,}939\) operations.

\subsection{A signed frame lowers the source degree}

For the source used in this endpoint, one may replace \(q^2W_0\) by
\(qW_0\). During completeness, the semantic construction gives
\(q=(BH)^x\), with \(x\ge2\), and every prime dividing the odd modulus
\(\mathcal M=Nq^x\) divides \(q\). Thus every such prime has square
dividing \(W=qW_0\). The proof of Theorem~\ref{thm:reuse} only needs
that local square divisibility, so the same modular construction works.
The positivity of the semantic \(q\) gives a positive completeness frame.

Soundness must also cover a negative value of the unclassified input \(q\).
The affine source row gives \(|q|\ge3\), so \(F=|W|\ge3W_0\).
The new \(A_*=L_0^2+2r^2\) is positive and at least \(L_0^2\).
When \(W=-F\), the frame equation becomes
\[
 L_0^2-F^3(F-2)r^2=1.
\]
This is the Pell equation at parameter \(F-1\), with ordinate \(F|r|\).
The congruence
\[
 \PellY_j(F-1)\equiv(-1)^{j-1}j\pmod F
\]
forces its positive classified index to be a multiple of \(F\).
As \(r\) is odd, it is nonzero, so that index is at least \(F\).
Consequently
\[
 A_*\ge L_0^2\ge(F-1)^{2F}>F^F\qquad(F\ge3).
\]
For \(W>0\) the retained positive-frame argument gives the required
bound. Every old bounded atom has absolute value below \(F\).
The source's only uses of \(W\) and its private frame wires are the frame
equations and its definition, checked by an occurrence audit.
The inherited signed power recognition therefore applies with \(F=|W|\).

The signed-frame row has degree three, rather than four.
Replaying the inherited integral coordinate changes, with the two private
center inputs deleted, gives an explicit equation in \(227\) integers,
the sum of \(175\) residual squares, of exact degree six.
Its collected polynomial has \(32{,}328\) monomials and height \(221{,}184\).
The archive also retains the positive-frame degree-eight alternative.
The \(227\)-variable intermediate is a proof source, not a global quartic
or low-dimensional record.

\subsection{Acyclic unit elimination}

Let \(S_0,\ldots,S_{174}\) be the literal residuals in that \(227\)-variable
source. A selected row has the form
\[
 S_r=\epsilon x_v+H,\qquad \epsilon\in\{1,-1\},
 \qquad H\in\Z[x_0,\ldots,\widehat{x_v},\ldots,x_{226}].
 \tag{Pivot}
\]
Set \(x_v=-\epsilon H\) and remove that row.
The dependence graph of the selected definitions is acyclic, so the
removed coordinates are restored as integer polynomials in the retained
ones, in a finite order. No division by a nonunit is performed.

For the principal endpoint the certificate contains \(144\) such
definitions, leaving \(83\) coordinates and \(31\) residuals.
If \(\Phi:\Z^{83}\to\Z^{227}\) is the resulting reconstruction, exact
polynomial substitution gives
\[
 S_r(\Phi(\boldsymbol x))=0
 \quad\text{for every removed row},\qquad
 D_{227}(\Phi(\boldsymbol x))=D_{83}(\boldsymbol x).
 \tag{Restore}
\]
Projection and reconstruction are inverse on the integer zero sets.
The signed-frame and witness-reuse changes earlier preserve existence;
the acyclic elimination itself gives this stronger zero-set bijection.
The method of unit elimination is inherited from the repository's prior
projection work. The new semantic source, selected dependency orientation,
literal certificates, and measured endpoints are supplied here.

\subsection{Exact degree and cost certificates}

The circuit gives an upper bound \(312\) on each retained residual degree.
A univariate specialization \(x_i=d_i t\), with explicitly supplied
integer \(d_i\), has a nonzero \(t^{312}\) coefficient in a retained
residual modulo \(1{,}000{,}003\). This proves that residual has exact
degree \(312\) over \(\Z\).
The sum of squares has degree \(624\): highest real homogeneous squares
cannot cancel. Independent finite-field nonzero partial derivatives
certify that every one of the \(83\) declared inputs actually occurs.

Exact propagation of coefficient-norm upper bounds through the circuit
gives the rigorous bound
\[
 \max |\text{expanded coefficient of }D_{83}|<2^{1568}.
\]
This is a bound, not an exact expanded height. A literal scalar bound
of \(64\) in the factored evaluator is a different measure. The complete expanded degree-\(624\)
polynomial was not collected; the circuit is its exact finite definition.


% BEGIN retained input: factored_evaluator.tex
\subsection{A factored evaluator for the same polynomial}
\label{sec:factored-evaluator}

The expanded source-coordinate presentation contains two long expressions
that share a quadratic form. Factoring these expressions changes neither
the source, the projection, nor the polynomial. It gives a second literal
evaluator with \textbf{1,789 live binary operations}, comprising 955
multiplications and 834 additions. The largest absolute literal is 64.
The canonical evaluator has 2,939 operations and largest literal 324.
Both define exactly the same polynomial in 83 unrestricted integers, of
degree 624, with 31 residual squares and expanded coefficient height below
$2^{1568}$.

Use the source coordinates $z_0,\ldots,z_{226}$ from the projection appendix,
and introduce the following polynomial abbreviations:
\begin{align*}
 H_*={}&z_8+z_9+4(z_{19}+z_{24}+z_{30}+z_{36}+z_{65})\\
       &+2(z_{22}+z_{27}+z_{29}+z_{32}+z_{34}
                    +z_{37}+z_{40}+z_{44}+z_{66}),\\
 q_*={}&3+6z_{220},\\
 \mathcal I={}&\{1,2,10,11,12,13,14,15,16,17,103,121,127,195\},\\
 \mathcal J={}&\{145,153,158,161,164,167,170,173,
                  177,180,183,187,190,214,215,216\},\\
 Q_*={}&z_0^2+(z_0+z_4)^2+\sum_{j\in\mathcal I}z_j^2,\\
 N_*={}&z_{193}^2+2z_{158}-z_{216}-\sum_{j\in\mathcal J}z_j^2.
\end{align*}
The long restoration of $z_{139}$ has the exact factorization
\begin{align}
 z_{139}={}&
 q_*\bigl[Q_*-2z_8+4z_{225}(z_{225}+1)+q_*^2+4\bigr]\notag\\
 &+2(1+3z_{220})(H_*^2+z_8^2+z_9^2)
       -2z_8z_9+z_{19}+N_*-3.\label{eq:factored-z139}
\end{align}
The corresponding retained residual is
\begin{equation}
 R_{28}=N_*-1-z_{19}-H_*^2+(z_8+z_9)^2.
 \label{eq:factored-R28}
\end{equation}
Here residual indexing is zero based: $R_{28}$ is retained source row 170.
These letters denote expression abbreviations, not additional unknowns.

The independent checker expands~\eqref{eq:factored-z139} and
\eqref{eq:factored-R28} in the free polynomial ring
$\mathbb Z[z_0,\ldots,z_{226}]$ and compares all 336 and 154 nonzero
coefficients, respectively, with the frozen definitions. It then checks
every restoration, every retained residual, and the sum-of-squares output
in both complete arithmetic DAGs. Thus the equality holds at every integer
tuple, not only on their zero sets. The original degree and occurrence
certificates also replay for the alternate evaluator. Exact propagation
of coefficient $\ell_1$ majorants again gives the strict bound $2^{1568}$;
this is distinct from the largest literal 64.

The alternate graph, formulas, byte-reproducible builder, independent
standard-library checker, and receipt are in
\texttt{direct/factored83/}. The canonical graph and certificates remain
in \texttt{direct/P83/}. No semantic theorem is changed by this evaluator
optimization.

% END retained input: factored_evaluator.tex


\begin{center}
\begin{tabular}{rrrr}
\toprule
Integer variables & Exact degree & Residual squares & Interpretation\\
\midrule
79 & 8,388,620 & 27 & Fewest coordinates in this family\\
81 & 2,048 & 29 & Smaller dimension, higher degree\\
82 & 764 & 30 & Intermediate point\\
\textbf{83} & \textbf{624} & \textbf{31} & \textbf{Principal direct endpoint}\\
87 & 268 & 35 & Lower-degree alternative\\
89 & 128 & 37 & Lower-degree alternative\\
91 & 80 & 39 & Lower-degree alternative\\
\bottomrule
\end{tabular}
\end{center}

Each row is an independently specified ordinary polynomial over unrestricted
integers, with its own literal circuit and exact degree certificate.
The \(79\)-variable equation improves on the \(87\)-variable coordinate
count in the inspected oriented-power family, at a slightly higher
degree than its \(8{,}388{,}612\) endpoint \cite{source177}.
The table makes the degree cost explicit. It is not a claim that an
equation in \(79\) variables is simpler than every equation in fewer
variables using the large native coefficient tuple.

% END retained input: direct_results.tex


\section{Verification, reproducibility, and remaining work}
The archive distinguishes three kinds of evidence:
\begin{enumerate}
\item \textbf{Literal arithmetic.} JSON circuits and residual lists specify
every polynomial, input, fixed coefficient, and operation.
The four native integers are supplied as bytes and by a finite reconstruction.
\item \textbf{Finite certificates.} Exact sparse algebra checks the center
identity, its norm defect, highest forms, and gate counts.
Separate code checks source edits, acyclic unit eliminations, reconstruction
identities, and the supplied collected coefficients.
Finite Pell examples exercise the lemma and a counterexample; they are
diagnostics, not the proof of the theorem.
\item \textbf{Unbounded reasoning.} The signed transport, index permutation,
cofinal completeness, and inherited refutation calibrations are ordinary
mathematical proofs. No finite search or successful test suite establishes
these universal statements by itself.
\end{enumerate}

Run \texttt{verify.py} for the release checks and
\texttt{build\_report.sh} to compile the report.
The single-file \texttt{report.tex} assembles the readable article and
coefficient specification. Electronic appendices contain the direct
polynomials, complete restoration data, source snapshots, and receipts.
The finite code does not search for an actual refutation or integer zero.

The results give concrete new equations and a reusable way to identify
two witness roles. Formalization of the new bridge end to end, and
further improvements under a specified simplicity objective, remain open.
No worldwide novelty priority is asserted.
The independent reviews were performed by separately assigned AI agents
and the primary assistant, not by external human peer reviewers.

\appendix
\section{The complete four-integer coefficient specification}\label{app:constants}
This is the retained directly indexed BFS943/radix67 tuple.
It makes \(P_{190}\) and \(P_{202}\) fully specified without an unspecified
machine, oracle coefficient, or consistency-dependent search.
Finite sums, floors, and the fixed seed below define four constants once
and for all. They do not add primitives to the polynomial in ten variables.

% BEGIN retained input: coefficient_recipe.tex
% Exact finite recipe extracted from the retained indexed_constants.tex.
% In this appendix K denotes the fixed port K_0, and all other symbols are constants.
\begingroup
\providecommand{\fl}[1]{\left\lfloor#1\right\rfloor}
\providecommand{\rg}[2]{\{#1,\ldots,#2\}}
In this appendix $K$ denotes the port $K_0$ in the new equations.
Every symbol introduced here is a fixed constant, independent of the ten
quantified variables. The formulas and 436-character seed are the retained
indexed BFS943 radix67 specification; they reconstruct the exact supplied
payloads, with $\beta=2^{590069}-1$.

\textbf{The four fixed integers.}
Every index, set, sum, product and floor below defines a constant, not
another unknown. Put
\[
\rho=2^{67},\quad B=\rho^{8807},\quad\beta=B-1,\quad
G_m(x)=\sum_{j=0}^{m-1}x^j.
\]
Define
\[
\begin{split}
\mathcal E={}&\rg0{942}\cup\rg{983}{992}\cup\rg{1003}{1011}\\
&\cup\rg{1013}{1031}\cup\{1063,1064,1095,2195,3296\}.
\end{split}
\]
Then
\begin{align*}
M_C={}&B-1-\sum_{e\in\mathcal E\setminus\{1\}}\rho^e-2\rho^{1063},\\
M_F={}&B-1+4+(2^{51}-2)(\rho^{4393}+\rho^{5495})\\
&+\rho^{1013}G_{10}(\rho)+\rho^{1065}G_9(\rho)
  +\rho^{2195}+\rho^{3296}\\
&+(2^{65}-1)\rho^{8792}G_{13}(\rho)+(2^{19}-1)\rho^{8805},\\
K={}&\rho^{30}+\rho^{62}+\rho^{1100}+\rho^{2201}+\rho^{8849}\\
&+(\rho^{4393}+\rho^{5495})\Phi(1/\rho),
\end{align*}
where
\[
\begin{split}
\Phi(x)={}&(1+2^{51})G_{943}(x)
 -2^{10}\sum_{i=0}^{942}v(t_i)x^i\\
&+\sum_{(b,e,m)\in\mathcal C}2^b x^e G_m(2x),\\
\mathcal C={}&\{(10,1023,8),(18,1013,9),(27,1003,8),\\
&(35,983,9),(44,1031,1),(45,1022,1),(46,1011,1),\\
&(47,992,1),(48,1095,1),(49,2195,1),(50,3296,1)\}.
\end{split}
\]
The rational intermediate in K cancels exactly; all four results are positive integers.

\textbf{The 943 indexed fixed tiles.}
Use the actions $u_j,v_j,d_j$ decoded below. Let
\[
\begin{gathered}
h(q,b)=2(q+1)+b,\quad
\mathcal R=\{u_j:d_j=1,\ 0\le j<430\},\\
\mathcal L=\{u_j:d_j=0,\ 0\le j<430\},\\
r_q=6+|\mathcal R\cap\rg0{q-1}|,\quad
l_q=205+|\mathcal L\cap\rg0{q-1}|.
\end{gathered}
\]
An interval ending at $-1$ is empty. The table defines $t_i$ for
$0\le i\le942$; all index ranges are fixed:
\[
\setlength{\arraycolsep}{1.4pt}\begin{array}{c|l}
\text{index }i&\text{tile }t_i=(w,n,e,s)\\\hline
0,1&(2,435,3,1),\ (5,0,0,0)\\
j+2\ (j=2q+b<430)&((1-d_j)l_{u_j},h(q,b),d_jr_{u_j},v_j)\\
432+b\ (b=0,1)&(0,h(215,b),0,h(215,b))\\
434+b\ (b=0,1)&(0,b,0,b)\\
424+2r_q+b\ (q\in\mathcal R,b=0,1)&(r_q,b,0,h(q,b))\\
424+2l_q+b\ (q\in\mathcal L,b=0,1)&(0,b,l_q,h(q,b))\\
924+j\ (0\le j<4)&(0,j+430\fl{j/2},0,438)\\
928&(2,434,2,0)\\
929&(3,436,4,4)\\
930&(4,436,4,0)\\
497+t\ (t=434,435,436)&(1,438,1,t)\\
\substack{934+3a+b-6\fl{a/4}\\a\in\{0,1,4\},\ b=0,1,2}&(a,437,b,437)
\end{array}
\]
For $t=(w,n,e,s)$ put
\[
v(t)=w+2^8n+2^{17}e+2^{25}s+127\,2^{34}.
\]

\textbf{A closed decoder for the 430 actions.}
Put $D_b(A,k)=\fl{A/b^k}\bmod b$, where
$a\bmod b=a-b\fl{a/b}$. For $0\le q<215$, $b=0,1$, $j=2q+b$:
\[
\begin{gathered}
N=\mathcal A\bmod2^{636},\quad
T_0=\fl{\mathcal A/2^{636}}\bmod2^{1675},\\
E_0=\fl{\mathcal A/2^{2311}},\\
\nu_j=D_2(N,j),\quad \sigma_j=\sum_{k=0}^{j-1}\nu_k,\\
m_q=D_2(N,430+q),\quad \tau_q=\sum_{k=0}^{q-1}m_k,\\
F_q=14(1-m_q)+m_qD_{16}(E_0,\tau_q),\\
u_j=\nu_j(1+\sigma_j)+(1-\nu_j)D_{216}(T_0,j-\sigma_j),\\
v_j=D_2(F_q,2b),\quad d_j=D_2(F_q,2b+1).
\end{gathered}
\]
Empty sums are zero. Here $(v_0,d_0,u_0)=(1,1,1)$; the constructor tile
uses this actual first target, including its state renaming.

\textbf{The integer $\mathcal A$.}
The alphabet, in order of digit values $0,\ldots,61$, is:\par
\noindent{\fontsize{7.7}{8.2}\selectfont\ttfamily 0123456789ABCDEFGHIJKLMNOPQRSTUVWXYZabcdefghijklmnopqrstuvwxyz}\par
Concatenate the following 436 characters as digits $e_0,\ldots,e_{435}$;
put $\mathcal A=\sum_{j=0}^{435}e_j62^{435-j}$.\par
\begingroup\fontsize{7.7}{8.2}\selectfont\ttfamily
BxEnhvy6NipzkWmROo0GsA3uK3twWLxscLq20UTumy5Cdu461LwD6g7RDNzF8Y\par
iYDFTpB8ehAypd7d6OsOL4APamnJeNXQypUezCmW9n4oXP6XfgmiWuAWZ8ZJ7U\par
OP7gu4kVdlGZqZHULQMar1dXNeJW4i7oNVI8usp8BlovjjZbwPCovZI328fRUq\par
4q92R81s5GjWQaQfGhkq9UYMK3rRMaBmsOIBkJqc2cj2VRIZ6mjzkBIOkW9SjZ\par
FzeA6J5u4j3sBL3SdgwXzM9mHvQivu8EJ1TBTpRhHDqZIbmYgwQ8VqTAFotl6F\par
We1YC554GKrVHsA8oUGKSA6F6k5KREBSD0Tgd7bFALY84LLbpw574BQyvdoHrl\par
0Oe35uwDf15DpT6VJWZBjSbHPpPAINYTte0eUxnTFTKBlGj8Am6sb1gLldh8WZ\par
8P\par
\endgroup

\endgroup

% END retained input: coefficient_recipe.tex


\section{Literal ten-variable circuit}
The following generated table prints the principal circuit in order.
Names denote local polynomial registers; the four fixed ports do not count
as inputs. Every row is one binary operation, and all \(118\) rows are live.

% BEGIN retained input: circuit.tex
\begingroup\fontsize{8.5}{10.5}\selectfont
\begin{longtable}{rll}
\toprule
Row & Register & Definition\\
\midrule\endhead
1 & \texttt{bJ} & $\mathtt{beta} \mathbin{\cdot} \mathtt{J}$ \\
2 & \texttt{q} & $\mathtt{bJ} + \mathtt{1}$ \\
3 & \texttt{q2} & $\mathtt{q} \mathbin{\cdot} \mathtt{q}$ \\
4 & \texttt{q4} & $\mathtt{q2} \mathbin{\cdot} \mathtt{q2}$ \\
5 & \texttt{q6} & $\mathtt{q4} \mathbin{\cdot} \mathtt{q2}$ \\
6 & \texttt{q7} & $\mathtt{q6} \mathbin{\cdot} \mathtt{q}$ \\
7 & \texttt{sixq7} & $\mathtt{6} \mathbin{\cdot} \mathtt{q7}$ \\
8 & \texttt{qY} & $\mathtt{q} \mathbin{\cdot} \mathtt{s}$ \\
9 & \texttt{qmF} & $\mathtt{q} - \mathtt{F}$ \\
10 & \texttt{C} & $\mathtt{qmF} - \mathtt{alpha}$ \\
11 & \texttt{X} & $\mathtt{w} \mathbin{\cdot} \mathtt{q}$ \\
12 & \texttt{k} & $\mathtt{eta} + \mathtt{zeta}$ \\
13 & \texttt{kY} & $\mathtt{k} \mathbin{\cdot} \mathtt{s}$ \\
14 & \texttt{c} & $\mathtt{kY} + \mathtt{eta}$ \\
15 & \texttt{XY} & $\mathtt{X} \mathbin{\cdot} \mathtt{qY}$ \\
16 & \texttt{v0} & $\mathtt{XY} \mathbin{\cdot} \mathtt{qY}$ \\
17 & \texttt{Am1} & $\mathtt{XY} + \mathtt{qY}$ \\
18 & \texttt{4Am1} & $\mathtt{4} \mathbin{\cdot} \mathtt{Am1}$ \\
19 & \texttt{Hp} & $\mathtt{4Am1} - \mathtt{1}$ \\
20 & \texttt{Ap1} & $\mathtt{Am1} + \mathtt{2}$ \\
21 & \texttt{V0k} & $\mathtt{v0} \mathbin{\cdot} \mathtt{k}$ \\
22 & \texttt{V0kk} & $\mathtt{V0k} + \mathtt{k}$ \\
23 & \texttt{V0k2} & $\mathtt{V0k} \mathbin{\cdot} \mathtt{V0kk}$ \\
24 & \texttt{Dtau} & $\mathtt{V0k2} + \mathtt{1}$ \\
25 & \texttt{G1} & $\mathtt{bJ} \mathbin{\cdot} \mathtt{qmF}$ \\
26 & \texttt{G} & $\mathtt{G1} + \mathtt{alpha}$ \\
27 & \texttt{q2m1} & $\mathtt{q2} - \mathtt{1}$ \\
28 & \texttt{Gq} & $\mathtt{G} \mathbin{\cdot} \mathtt{q2m1}$ \\
29 & \texttt{qMF} & $\mathtt{q} \mathbin{\cdot} \mathtt{MF}$ \\
30 & \texttt{MCqMF} & $\mathtt{MC} + \mathtt{qMF}$ \\
31 & \texttt{RJ} & $\mathtt{MCqMF} \mathbin{\cdot} \mathtt{J}$ \\
32 & \texttt{R} & $\mathtt{Gq} + \mathtt{RJ}$ \\
33 & \texttt{Rp1} & $\mathtt{R} + \mathtt{1}$ \\
34 & \texttt{e} & $\mathtt{k} - \mathtt{Rp1}$ \\
35 & \texttt{ksq} & $\mathtt{k} \mathbin{\cdot} \mathtt{k}$ \\
36 & \texttt{ksqm4} & $\mathtt{ksq} - \mathtt{4}$ \\
37 & \texttt{inverse\_cubic} & $\mathtt{V0k} \mathbin{\cdot} \mathtt{ksqm4}$ \\
38 & \texttt{q4s} & $\mathtt{q4} \mathbin{\cdot} \mathtt{s}$ \\
39 & \texttt{e\_sigma} & $\mathtt{e} - \mathtt{q4s}$ \\
40 & \texttt{six\_e\_sigma} & $\mathtt{6} \mathbin{\cdot} \mathtt{e\_sigma}$ \\
41 & \texttt{eclear} & $\mathtt{six\_e\_sigma} - \mathtt{inverse\_cubic}$ \\
42 & \texttt{K0w} & $\mathtt{K0} + \mathtt{w}$ \\
43 & \texttt{K0wC} & $\mathtt{K0w} \mathbin{\cdot} \mathtt{C}$ \\
44 & \texttt{Bt} & $\mathtt{K0wC} - \mathtt{F}$ \\
45 & \texttt{c2} & $\mathtt{c} \mathbin{\cdot} \mathtt{c}$ \\
46 & \texttt{U} & $\mathtt{i} \mathbin{\cdot} \mathtt{c2}$ \\
47 & \texttt{UE} & $\mathtt{U} \mathbin{\cdot} \mathtt{Dtau}$ \\
48 & \texttt{Delta} & $\mathtt{Am1} \mathbin{\cdot} \mathtt{Ap1}$ \\
49 & \texttt{SE} & $\mathtt{Delta} \mathbin{\cdot} \mathtt{UE}$ \\
50 & \texttt{DpI2} & $\mathtt{SE} \mathbin{\cdot} \mathtt{U}$ \\
51 & \texttt{Df} & $\mathtt{DpI2} + \mathtt{1}$ \\
52 & \texttt{aF0} & $\mathtt{DpI2} \mathbin{\cdot} \mathtt{Df}$ \\
53 & \texttt{fouraF0} & $\mathtt{4} \mathbin{\cdot} \mathtt{aF0}$ \\
54 & \texttt{trace\_f} & $\mathtt{DpI2} + \mathtt{Df}$ \\
55 & \texttt{Q} & $\mathtt{Delta} \mathbin{\cdot} \mathtt{fouraF0}$ \\
56 & \texttt{cT} & $\mathtt{c} \mathbin{\cdot} \mathtt{T}$ \\
57 & \texttt{cTR} & $\mathtt{cT} - \mathtt{R}$ \\
58 & \texttt{DfcTR} & $\mathtt{trace\_f} \mathbin{\cdot} \mathtt{cTR}$ \\
59 & \texttt{V} & $\mathtt{DfcTR} - \mathtt{c}$ \\
60 & \texttt{V2} & $\mathtt{V} \mathbin{\cdot} \mathtt{V}$ \\
61 & \texttt{P1mu2} & $\mathtt{bJ} \mathbin{\cdot} \mathtt{sixq7}$ \\
62 & \texttt{new\_A} & $\mathtt{Am1} + \mathtt{1}$ \\
63 & \texttt{new\_kZ} & $\mathtt{k} \mathbin{\cdot} \mathtt{qY}$ \\
64 & \texttt{new\_Zp1} & $\mathtt{qY} + \mathtt{1}$ \\
65 & \texttt{new\_kZZp1} & $\mathtt{new\_kZ} \mathbin{\cdot} \mathtt{new\_Zp1}$ \\
66 & \texttt{new\_Sinner} & $\mathtt{V0kk} - \mathtt{new\_kZZp1}$ \\
67 & \texttt{new\_S} & $\mathtt{new\_kZ} \mathbin{\cdot} \mathtt{new\_Sinner}$ \\
68 & \texttt{new\_e2} & $\mathtt{new\_kZ} \mathbin{\cdot} \mathtt{new\_kZ}$ \\
69 & \texttt{new\_U2} & $\mathtt{U} \mathbin{\cdot} \mathtt{U}$ \\
70 & \texttt{new\_theta} & $\mathtt{new\_U2} \mathbin{\cdot} \mathtt{cTR}$ \\
71 & \texttt{new\_threeE} & $\mathtt{3} \mathbin{\cdot} \mathtt{Dtau}$ \\
72 & \texttt{new\_3Ee2} & $\mathtt{new\_threeE} + \mathtt{new\_e2}$ \\
73 & \texttt{new\_theta3Ee2} & $\mathtt{new\_theta} \mathbin{\cdot} \mathtt{new\_3Ee2}$ \\
74 & \texttt{new\_bmc} & $\mathtt{cTR} - \mathtt{c}$ \\
75 & \texttt{new\_Rc} & $\mathtt{new\_theta3Ee2} - \mathtt{new\_bmc}$ \\
76 & \texttt{new\_E2} & $\mathtt{Dtau} \mathbin{\cdot} \mathtt{Dtau}$ \\
77 & \texttt{new\_fourE2} & $\mathtt{4} \mathbin{\cdot} \mathtt{new\_E2}$ \\
78 & \texttt{new\_S2} & $\mathtt{new\_S} \mathbin{\cdot} \mathtt{new\_S}$ \\
79 & \texttt{new\_threeS2} & $\mathtt{3} \mathbin{\cdot} \mathtt{new\_S2}$ \\
80 & \texttt{new\_e22} & $\mathtt{new\_e2} \mathbin{\cdot} \mathtt{new\_e2}$ \\
81 & \texttt{new\_sc1} & $\mathtt{new\_fourE2} + \mathtt{new\_threeS2}$ \\
82 & \texttt{new\_Sc} & $\mathtt{new\_sc1} + \mathtt{new\_e22}$ \\
83 & \texttt{new\_bB} & $\mathtt{cTR} \mathbin{\cdot} \mathtt{Bt}$ \\
84 & \texttt{new\_ScShift} & $\mathtt{new\_Sc} - \mathtt{new\_bB}$ \\
85 & \texttt{new\_EV} & $\mathtt{Dtau} \mathbin{\cdot} \mathtt{V}$ \\
86 & \texttt{new\_fourEV} & $\mathtt{4} \mathbin{\cdot} \mathtt{new\_EV}$ \\
87 & \texttt{new\_thetaSc} & $\mathtt{new\_theta} \mathbin{\cdot} \mathtt{new\_ScShift}$ \\
88 & \texttt{new\_Vdiff} & $\mathtt{new\_fourEV} - \mathtt{new\_thetaSc}$ \\
89 & \texttt{new\_Apart} & $\mathtt{new\_A} \mathbin{\cdot} \mathtt{new\_Vdiff}$ \\
90 & \texttt{new\_SRc} & $\mathtt{new\_S} \mathbin{\cdot} \mathtt{new\_Rc}$ \\
91 & \texttt{new\_twoSRc} & $\mathtt{2} \mathbin{\cdot} \mathtt{new\_SRc}$ \\
92 & \texttt{new\_omegainner} & $\mathtt{new\_Apart} + \mathtt{new\_twoSRc}$ \\
93 & \texttt{new\_Omega} & $\mathtt{U} \mathbin{\cdot} \mathtt{new\_omegainner}$ \\
94 & \texttt{center} & $\mathtt{n} - \mathtt{new\_Omega}$ \\
95 & \texttt{Hn} & $\mathtt{Hp} \mathbin{\cdot} \mathtt{center}$ \\
96 & \texttt{threec} & $\mathtt{3} \mathbin{\cdot} \mathtt{c}$ \\
97 & \texttt{twoX} & $\mathtt{2} \mathbin{\cdot} \mathtt{X}$ \\
98 & \texttt{gap} & $\mathtt{threec} - \mathtt{twoX}$ \\
99 & \texttt{Kpos0} & $\mathtt{X} \mathbin{\cdot} \mathtt{gap}$ \\
100 & \texttt{Kpos} & $\mathtt{Kpos0} + \mathtt{2}$ \\
101 & \texttt{lin} & $\mathtt{Hn} + \mathtt{Kpos}$ \\
102 & \texttt{Zpart1} & $\mathtt{P1mu2} \mathbin{\cdot} \mathtt{lin}$ \\
103 & \texttt{mu2inner} & $\mathtt{sixq7} \mathbin{\cdot} \mathtt{Bt}$ \\
104 & \texttt{P1e} & $\mathtt{bJ} \mathbin{\cdot} \mathtt{eclear}$ \\
105 & \texttt{brk} & $\mathtt{mu2inner} + \mathtt{P1e}$ \\
106 & \texttt{Zpart2} & $\mathtt{Hp} \mathbin{\cdot} \mathtt{brk}$ \\
107 & \texttt{Znum} & $\mathtt{Zpart1} - \mathtt{Zpart2}$ \\
108 & \texttt{Mtot} & $\mathtt{P1mu2} \mathbin{\cdot} \mathtt{Hp}$ \\
109 & \texttt{M2} & $\mathtt{Mtot} \mathbin{\cdot} \mathtt{Mtot}$ \\
110 & \texttt{Dg} & $\mathtt{Dtau} \mathbin{\cdot} \mathtt{Df}$ \\
111 & \texttt{new\_scaledG} & $\mathtt{16} \mathbin{\cdot} \mathtt{Dg}$ \\
112 & \texttt{Z2} & $\mathtt{Znum} \mathbin{\cdot} \mathtt{Znum}$ \\
113 & \texttt{RG} & $\mathtt{M2} \mathbin{\cdot} \mathtt{new\_scaledG}$ \\
114 & \texttt{RGV2} & $\mathtt{RG} \mathbin{\cdot} \mathtt{V2}$ \\
115 & \texttt{inside} & $\mathtt{Z2} - \mathtt{RGV2}$ \\
116 & \texttt{Qinside} & $\mathtt{Q} \mathbin{\cdot} \mathtt{inside}$ \\
117 & \texttt{correction} & $\mathtt{RG} - \mathtt{Z2}$ \\
118 & \texttt{N} & $\mathtt{Qinside} + \mathtt{correction}$ \\
\bottomrule\end{longtable}
\endgroup

% END retained input: circuit.tex



% BEGIN retained input: P83_definition.tex
\section{Complete finite definition of the direct polynomial}\label{app:direct83-definition}
The only independent unknowns in this appendix are $x_0,\ldots,x_{82}\in\mathbb Z$. Every $z_i$ below is an abbreviation for an integer polynomial in these unknowns. First make the free-coordinate assignments, then evaluate the acyclic definitions in their printed order. Finally evaluate the 31 residuals and add their squares. This is a complete ordinary polynomial definition using fixed integer coefficients, addition, and multiplication.
\begingroup\small\allowdisplaybreaks
\subsection*{Free-coordinate assignments}
\begin{align*}
z_{5}&=x_{0},\qquad z_{6}&=x_{1},\qquad z_{8}&=x_{2},\qquad z_{9}&=x_{3}\\
z_{10}&=x_{4},\qquad z_{11}&=x_{5},\qquad z_{12}&=x_{6},\qquad z_{13}&=x_{7}\\
z_{14}&=x_{8},\qquad z_{15}&=x_{9},\qquad z_{18}&=x_{10},\qquad z_{48}&=x_{11}\\
z_{53}&=x_{12},\qquad z_{62}&=x_{13},\qquad z_{72}&=x_{14},\qquad z_{101}&=x_{15}\\
z_{104}&=x_{16},\qquad z_{109}&=x_{17},\qquad z_{111}&=x_{18},\qquad z_{112}&=x_{19}\\
z_{115}&=x_{20},\qquad z_{118}&=x_{21},\qquad z_{121}&=x_{22},\qquad z_{122}&=x_{23}\\
z_{124}&=x_{24},\qquad z_{125}&=x_{25},\qquad z_{128}&=x_{26},\qquad z_{131}&=x_{27}\\
z_{132}&=x_{28},\qquad z_{134}&=x_{29},\qquad z_{136}&=x_{30},\qquad z_{137}&=x_{31}\\
z_{141}&=x_{32},\qquad z_{144}&=x_{33},\qquad z_{145}&=x_{34},\qquad z_{147}&=x_{35}\\
z_{149}&=x_{36},\qquad z_{151}&=x_{37},\qquad z_{152}&=x_{38},\qquad z_{153}&=x_{39}\\
z_{156}&=x_{40},\qquad z_{157}&=x_{41},\qquad z_{158}&=x_{42},\qquad z_{161}&=x_{43}\\
z_{163}&=x_{44},\qquad z_{164}&=x_{45},\qquad z_{166}&=x_{46},\qquad z_{167}&=x_{47}\\
z_{169}&=x_{48},\qquad z_{170}&=x_{49},\qquad z_{173}&=x_{50},\qquad z_{175}&=x_{51}\\
z_{177}&=x_{52},\qquad z_{179}&=x_{53},\qquad z_{180}&=x_{54},\qquad z_{183}&=x_{55}\\
z_{185}&=x_{56},\qquad z_{187}&=x_{57},\qquad z_{189}&=x_{58},\qquad z_{190}&=x_{59}\\
z_{193}&=x_{60},\qquad z_{196}&=x_{61},\qquad z_{197}&=x_{62},\qquad z_{202}&=x_{63}\\
z_{203}&=x_{64},\qquad z_{204}&=x_{65},\qquad z_{205}&=x_{66},\qquad z_{206}&=x_{67}\\
z_{210}&=x_{68},\qquad z_{211}&=x_{69},\qquad z_{212}&=x_{70},\qquad z_{213}&=x_{71}\\
z_{214}&=x_{72},\qquad z_{215}&=x_{73},\qquad z_{216}&=x_{74},\qquad z_{217}&=x_{75}\\
z_{218}&=x_{76},\qquad z_{219}&=x_{77},\qquad z_{220}&=x_{78},\qquad z_{222}&=x_{79}\\
z_{223}&=x_{80},\qquad z_{225}&=x_{81},\qquad z_{226}&=x_{82}
\end{align*}
\subsection*{Acyclic unit definitions}
\begin{align*}
z_{0}&=16+z_{196}^{2}\\[3pt]
z_{209}&=z_{193}^{2}\\[3pt]
z_{198}&=2z_{203} z_{209}+z_{209}\\[3pt]
z_{201}&=4z_{198}+8z_{198} z_{202}\\[3pt]
z_{140}&=-6-z_{197} z_{204}+z_{201}\\[3pt]
z_{146}&=z_{140} z_{145}+7z_{145}\\[3pt]
z_{4}&=1-32z_{140} z_{147}+z_{144}+16z_{145}-z_{146}+33z_{147}\\[3pt]
z_{3}&=z_{4}+2z_{4} z_{225}
\end{align*}
\begin{align*}
z_{221}&=1+4z_{222}+4z_{222}^{2}\\[3pt]
z_{2}&=1+z_{3} z_{221}\\[3pt]
z_{7}&=-z_{0}+z_{0} z_{2}+z_{2}\\[3pt]
z_{1}&=1+z_{6} z_{7}\\[3pt]
z_{16}&=-1-z_{4}+z_{18}+2z_{18} z_{225}\\[3pt]
z_{148}&=12+2z_{140}+4z_{140} z_{225}+24z_{225}-4z_{225}^{2}\\[3pt]
z_{224}&=3+2z_{0}+2z_{4}-2z_{8}+z_{8}^{2}+z_{9}^{2}+z_{10}^{2}+z_{11}^{2}+z_{12}^{2}+z_{13}^{2}+z_{14}^{2}+z_{15}^{2}\\*
&\quad{}+z_{16}^{2}\\[3pt]
z_{195}&=z_{224}+2z_{224} z_{225}
\end{align*}
\begin{align*}
z_{150}&=-1+2z_{140} z_{195}+14z_{195}-z_{195}^{2}\\[3pt]
z_{184}&=-1+z_{140} z_{183}+7z_{183}\\[3pt]
z_{182}&=3+z_{150} z_{151}-z_{183} z_{195}+z_{184}+6z_{220}\\[3pt]
z_{17}&=-z_{148} z_{149}+z_{182}+z_{183}+2z_{183} z_{225}-z_{184}\\[3pt]
z_{19}&=1+z_{8}^{2}+2z_{8} z_{9}+z_{9}^{2}\\[3pt]
z_{20}&=10z_{8}-z_{8}^{2}+6z_{8} z_{9}+8z_{8} z_{19}+20z_{9}+11z_{9}^{2}+16z_{9} z_{19}+15z_{19}+4z_{19}^{2}\\[3pt]
z_{21}&=2-2z_{8}+z_{8}^{2}+4z_{8} z_{20}+4z_{20}+4z_{20}^{2}\\[3pt]
z_{22}&=-1-6z_{8}+36z_{8}^{2}+24z_{8} z_{21}-2z_{21}+4z_{21}^{2}
\end{align*}
\begin{align*}
z_{23}&=2+z_{9}+z_{9}^{2}+2z_{9} z_{10}+z_{10}+z_{10}^{2}\\[3pt]
z_{46}&=1+8z_{10}+4z_{10}^{2}-8z_{10} z_{23}-8z_{23}+4z_{23}^{2}\\[3pt]
z_{24}&=8-7z_{8}+z_{8}^{2}-4z_{8} z_{10}+4z_{8} z_{23}+4z_{10}-4z_{23}+z_{46}\\[3pt]
z_{25}&=1-z_{8}+z_{8}^{2}+2z_{8} z_{10}-z_{10}+z_{10}^{2}\\[3pt]
z_{31}&=-19+25z_{8}+7z_{8}^{2}+8z_{8} z_{9}+30z_{8} z_{10}+16z_{8} z_{25}-3z_{9}-5z_{9}^{2}+4z_{9} z_{25}+27z_{10}\\*
&\quad{}+21z_{10}^{2}+24z_{10} z_{25}+6z_{23}+17z_{25}+4z_{25}^{2}\\[3pt]
z_{26}&=-23-16z_{8}+8z_{8} z_{9}+8z_{8} z_{10}+16z_{8} z_{19}+8z_{8} z_{25}-19z_{9}-9z_{9}^{2}+36z_{9} z_{10}\\*
&\quad{}+12z_{9} z_{25}-16z_{10}-8z_{10}^{2}+24z_{10} z_{19}+8z_{19} z_{25}+z_{20}+z_{31}\\[3pt]
z_{27}&=4z_{24}^{2}+8z_{24} z_{26}+4z_{26}^{2}\\[3pt]
z_{28}&=2+8z_{8}-8z_{8} z_{23}-8z_{8} z_{25}+8z_{9}-8z_{9} z_{23}-8z_{9} z_{25}+4z_{19}-8z_{23}+4z_{23}^{2}\\*
&\quad{}+8z_{23} z_{25}-8z_{25}+4z_{25}^{2}
\end{align*}
\begin{align*}
z_{29}&=-6-4z_{8}-32z_{8} z_{9}-16z_{8} z_{19}-16z_{8} z_{25}-8z_{8} z_{28}-24z_{9}+32z_{9} z_{25}+16z_{9} z_{28}\\*
&\quad{}-12z_{19}+16z_{19} z_{25}+8z_{19} z_{28}+z_{20}-8z_{25}+16z_{25}^{2}+16z_{25} z_{28}-4z_{28}+4z_{28}^{2}\\[3pt]
z_{30}&=-22+5z_{8}+z_{8}^{2}+8z_{8} z_{9}+12z_{8} z_{10}+4z_{8} z_{23}+8z_{9}+16z_{9} z_{10}+16z_{9} z_{23}+20z_{10}\\*
&\quad{}+20z_{10}^{2}+24z_{10} z_{23}+28z_{23}+4z_{23}^{2}\\[3pt]
z_{32}&=-2+4z_{30}^{2}+8z_{30} z_{31}+4z_{31}^{2}\\[3pt]
z_{33}&=19z_{8}-8z_{8}^{2}-8z_{8} z_{19}-8z_{8} z_{23}+46z_{9}+24z_{9}^{2}-24z_{9} z_{19}-24z_{9} z_{23}+2z_{10}\\*
&\quad{}-4z_{19}+4z_{19}^{2}+8z_{19} z_{23}-14z_{23}+4z_{23}^{2}+z_{25}\\[3pt]
z_{34}&=3+6z_{8}+4z_{10}+2z_{25}+4z_{33}+4z_{33}^{2}\\[3pt]
z_{38}&=6+2z_{8}+4z_{9}+14z_{19}+16z_{19}^{2}\\[3pt]
z_{35}&=-8-42z_{8}-7z_{8}^{2}-28z_{8} z_{20}-48z_{9}+32z_{9}^{2}+32z_{9} z_{19}-64z_{9} z_{20}+40z_{19}^{2}\\*
&\quad{}-32z_{19} z_{20}-14z_{20}+4z_{20}^{2}-2z_{38}\\[3pt]
z_{36}&=-1-6z_{8}+36z_{8}^{2}+24z_{8} z_{35}-2z_{35}+4z_{35}^{2}
\end{align*}
\begin{align*}
z_{37}&=z_{9}^{2}+4z_{9} z_{36}+4z_{36}^{2}\\[3pt]
z_{39}&=-24-21z_{9}+21z_{9}^{2}+18z_{9} z_{10}+24z_{9} z_{23}+24z_{9} z_{38}-3z_{10}+9z_{10}^{2}+16z_{10} z_{38}\\*
&\quad{}-3z_{23}+12z_{23}^{2}+8z_{23} z_{38}-4z_{38}+4z_{38}^{2}-2z_{46}\\[3pt]
z_{40}&=-16+10z_{8}+20z_{8}^{2}+16z_{8} z_{10}+24z_{8} z_{25}+24z_{8} z_{39}+12z_{10}+16z_{10} z_{25}\\*
&\quad{}+16z_{10} z_{39}+14z_{25}+4z_{25}^{2}+8z_{25} z_{39}-4z_{39}+4z_{39}^{2}\\[3pt]
z_{41}&=z_{10}^{2}+2z_{10} z_{11}+z_{11}^{2}\\[3pt]
z_{42}&=-2-z_{8}^{2}+4z_{8} z_{41}-2z_{10}+2z_{10}^{2}+8z_{10} z_{41}+2z_{25}-4z_{41}+4z_{41}^{2}\\[3pt]
z_{43}&=31+16z_{10}+16z_{10} z_{23}-16z_{10} z_{25}+8z_{23}-12z_{23}^{2}+8z_{23} z_{25}-24z_{25}+4z_{25}^{2}+4z_{46}\\[3pt]
z_{44}&=4z_{10}+16z_{25}^{2}+32z_{25} z_{41}+16z_{25} z_{42}+16z_{25} z_{43}+4z_{41}+16z_{41}^{2}+16z_{41} z_{42}\\*
&\quad{}+16z_{41} z_{43}+2z_{42}+4z_{42}^{2}+8z_{42} z_{43}+4z_{43}^{2}\\[3pt]
z_{45}&=16z_{8}^{2}-8z_{8} z_{10}-32z_{8} z_{25}-16z_{8} z_{28}-z_{10}+z_{10}^{2}+8z_{10} z_{25}+4z_{10} z_{28}\\*
&\quad{}+16z_{25}^{2}+16z_{25} z_{28}+4z_{28}^{2}
\end{align*}
\begin{align*}
z_{47}&=1+9z_{10}^{2}-24z_{10} z_{23}-12z_{10} z_{46}+16z_{23}^{2}+16z_{23} z_{46}+4z_{46}^{2}\\[3pt]
z_{49}&=1+4z_{48}+4z_{48}^{2}\\[3pt]
z_{50}&=1+z_{10}^{2}+8z_{10} z_{48}+4z_{10} z_{49}+16z_{48}^{2}+16z_{48} z_{49}+4z_{49}^{2}\\[3pt]
z_{51}&=-2-z_{9}+z_{9}^{2}+4z_{9} z_{10}+4z_{9} z_{50}+4z_{10}^{2}+8z_{10} z_{50}+4z_{50}^{2}\\[3pt]
z_{52}&=4z_{45}^{2}+8z_{45} z_{51}+4z_{51}^{2}\\[3pt]
z_{54}&=1+4z_{53}+4z_{53}^{2}\\[3pt]
z_{55}&=z_{10}^{2}+8z_{10} z_{53}+4z_{10} z_{54}+16z_{53}^{2}+16z_{53} z_{54}+4z_{54}^{2}\\[3pt]
z_{56}&=z_{8}+z_{8}^{2}+4z_{8} z_{10}+4z_{8} z_{55}+4z_{10}+4z_{10}^{2}+8z_{10} z_{55}+4z_{55}+4z_{55}^{2}
\end{align*}
\begin{align*}
z_{57}&=10+12z_{56}+4z_{56}^{2}\\[3pt]
z_{58}&=-2-4z_{10}-16z_{10} z_{56}-8z_{10} z_{57}+4z_{23}+16z_{23} z_{56}+8z_{23} z_{57}+z_{46}-8z_{56}+16z_{56}^{2}\\*
&\quad{}+16z_{56} z_{57}-4z_{57}+4z_{57}^{2}\\[3pt]
z_{59}&=-z_{9}+z_{9}^{2}-8z_{9} z_{10}+8z_{9} z_{23}+4z_{9} z_{58}+16z_{10}^{2}-32z_{10} z_{23}-16z_{10} z_{58}\\*
&\quad{}+16z_{23}^{2}+16z_{23} z_{58}+4z_{58}^{2}\\[3pt]
z_{60}&=10+12z_{51}+4z_{51}^{2}\\[3pt]
z_{61}&=-1+16z_{51}^{2}+16z_{51} z_{59}+16z_{51} z_{60}+4z_{59}^{2}+8z_{59} z_{60}+4z_{60}^{2}\\[3pt]
z_{76}&=z_{8} z_{9}-z_{8} z_{10}-z_{9}-z_{9} z_{10}+z_{10}+z_{10}^{2}\\[3pt]
z_{63}&=6z_{8}+3z_{8}^{2}+18z_{8} z_{9}+16z_{9}+27z_{9}^{2}-4z_{9} z_{10}-3z_{10}-3z_{10}^{2}+4z_{10} z_{19}\\*
&\quad{}+19z_{19}+8z_{19}^{2}-z_{20}+4z_{76}\\[3pt]
z_{64}&=-1+8z_{8}+4z_{8} z_{10}-8z_{8} z_{63}+16z_{9}+8z_{9} z_{10}-16z_{9} z_{63}-z_{10}-z_{10}^{2}-4z_{10} z_{19}\\*
&\quad{}-8z_{19}+8z_{19} z_{63}-3z_{63}+4z_{63}^{2}
\end{align*}
\begin{align*}
z_{65}&=2+12z_{8}+16z_{8} z_{10}-16z_{8} z_{62}-16z_{8} z_{64}+24z_{9}+32z_{9} z_{10}-32z_{9} z_{62}-32z_{9} z_{64}\\*
&\quad{}-4z_{10}-4z_{10}^{2}-16z_{10} z_{19}-12z_{19}+16z_{19} z_{62}+16z_{19} z_{64}+4z_{62}^{2}+8z_{62} z_{64}\\*
&\quad{}+4z_{63}+2z_{64}+4z_{64}^{2}\\[3pt]
z_{66}&=2+4z_{65}+4z_{65}^{2}\\[3pt]
z_{67}&=z_{13}+2z_{13} z_{225}\\[3pt]
z_{68}&=z_{12}-z_{13}-2z_{22} z_{67}\\[3pt]
z_{69}&=z_{12}-z_{13}-4z_{24} z_{67}-2z_{27} z_{67}\\[3pt]
z_{75}&=z_{12}-z_{13}-2z_{44} z_{67}-4z_{67}\\[3pt]
z_{70}&=-2z_{29} z_{67}+2z_{44} z_{67}-2z_{67}+z_{75}\\[3pt]
z_{71}&=z_{12}-z_{13}-4z_{30} z_{67}-2z_{32} z_{67}
\end{align*}
\begin{align*}
z_{73}&=z_{12}-z_{13}-4z_{36} z_{67}-2z_{37} z_{67}\\[3pt]
z_{74}&=z_{12}-z_{13}-2z_{40} z_{67}\\[3pt]
z_{77}&=-1-2z_{11}-z_{11}^{2}+z_{76}^{2}\\[3pt]
z_{78}&=z_{77}^{2}\\[3pt]
z_{79}&=4z_{65} z_{67}+2z_{66} z_{67}-2z_{67}\\[3pt]
z_{80}&=z_{12}^{2}-2z_{12} z_{13}-2z_{12} z_{79}+z_{13}^{2}+2z_{13} z_{79}+z_{79}^{2}\\[3pt]
z_{81}&=1+z_{9}+2z_{9} z_{225}\\[3pt]
z_{82}&=z_{8}+2z_{8} z_{225}-2z_{225}
\end{align*}
\begin{align*}
z_{83}&=-1-2z_{8}-4z_{8} z_{225}-4z_{9}-8z_{9} z_{225}+2z_{19}+4z_{19} z_{225}-4z_{225}\\[3pt]
z_{84}&=z_{12} z_{82}\\[3pt]
z_{85}&=z_{13} z_{81}-z_{83} z_{84}\\[3pt]
z_{86}&=-z_{12} z_{81}-6z_{13}+8z_{13} z_{81}+4z_{13} z_{82}+z_{13} z_{83}+12z_{13} z_{225}\\[3pt]
z_{87}&=z_{12} z_{68}-z_{13} z_{68}\\[3pt]
z_{88}&=z_{69} z_{87}\\[3pt]
z_{89}&=z_{70} z_{88}\\[3pt]
z_{90}&=z_{71} z_{89}
\end{align*}
\begin{align*}
z_{91}&=z_{72} z_{90}\\[3pt]
z_{92}&=z_{73} z_{91}\\[3pt]
z_{93}&=z_{74} z_{92}\\[3pt]
z_{94}&=z_{75} z_{93}\\[3pt]
z_{95}&=z_{85} z_{94}\\[3pt]
z_{96}&=z_{86} z_{95}\\[3pt]
z_{97}&=z_{96}^{2}\\[3pt]
z_{98}&=z_{2}+z_{2} z_{16}-z_{16}
\end{align*}
\begin{align*}
z_{99}&=z_{14} z_{98}\\[3pt]
z_{100}&=z_{2} z_{15}-z_{15}+z_{15} z_{98}\\[3pt]
z_{102}&=1-z_{140} z_{153}+z_{152}-3z_{153}+6z_{153} z_{220}\\[3pt]
z_{127}&=z_{121} z_{213}\\[3pt]
z_{130}&=z_{17} z_{127}\\[3pt]
z_{155}&=-7+z_{130} z_{226}-2z_{140}+12z_{220}\\[3pt]
z_{103}&=1+z_{144}+3z_{145}+6z_{145} z_{220}-z_{146}-z_{155} z_{156}\\[3pt]
z_{106}&=3z_{104}+6z_{104} z_{220}
\end{align*}
\begin{align*}
z_{105}&=z_{102} z_{211}-z_{106} z_{127}-z_{127}-2z_{127} z_{225}\\[3pt]
z_{107}&=4+z_{140}-z_{140} z_{158}+z_{157}-3z_{158}+6z_{158} z_{220}-6z_{220}\\[3pt]
z_{108}&=3z_{103}+6z_{103} z_{220}\\[3pt]
z_{110}&=z_{107} z_{212}-z_{109} z_{130}\\[3pt]
z_{113}&=z_{111} z_{212}-z_{112} z_{130}\\[3pt]
z_{114}&=1-z_{140} z_{164}+z_{163}-3z_{164}+6z_{164} z_{220}\\[3pt]
z_{116}&=z_{114} z_{212}-z_{115} z_{130}\\[3pt]
z_{117}&=1-z_{140} z_{167}+z_{166}-3z_{167}+6z_{167} z_{220}
\end{align*}
\begin{align*}
z_{119}&=z_{117} z_{212}-z_{118} z_{130}\\[3pt]
z_{120}&=1-z_{140} z_{170}+z_{169}-3z_{170}+6z_{170} z_{220}\\[3pt]
z_{123}&=z_{120} z_{213}-z_{122} z_{130}\\[3pt]
z_{126}&=z_{124} z_{213}-z_{125} z_{130}\\[3pt]
z_{129}&=1-z_{128} z_{130}-z_{140} z_{180}+z_{179}-3z_{180}+6z_{180} z_{220}\\[3pt]
z_{133}&=1-z_{140} z_{190}+z_{189}-3z_{190}+6z_{190} z_{220}\\[3pt]
z_{135}&=-z_{130} z_{134}+z_{133} z_{211}\\[3pt]
z_{138}&=-z_{130} z_{137}+z_{136} z_{212}
\end{align*}
\begin{align*}
z_{139}&=36+6z_{0}^{2}+12z_{0}^{2} z_{220}+6z_{0} z_{4}+12z_{0} z_{4} z_{220}+3z_{1}^{2}+6z_{1}^{2} z_{220}\\
&\quad{}+3z_{2}^{2}+6z_{2}^{2} z_{220}+3z_{4}^{2}+6z_{4}^{2} z_{220}-6z_{8}+4z_{8}^{2}+12z_{8}^{2} z_{220}\\
&\quad{}+2z_{8} z_{9}+12z_{8} z_{9} z_{220}+16z_{8} z_{19}+48z_{8} z_{19} z_{220}+8z_{8} z_{22}\\
&\quad{}+24z_{8} z_{22} z_{220}+16z_{8} z_{24}+48z_{8} z_{24} z_{220}+8z_{8} z_{27}+24z_{8} z_{27} z_{220}\\
&\quad{}+8z_{8} z_{29}+24z_{8} z_{29} z_{220}+16z_{8} z_{30}+48z_{8} z_{30} z_{220}+8z_{8} z_{32}\\
&\quad{}+24z_{8} z_{32} z_{220}+8z_{8} z_{34}+24z_{8} z_{34} z_{220}+16z_{8} z_{36}+48z_{8} z_{36} z_{220}\\
&\quad{}+8z_{8} z_{37}+24z_{8} z_{37} z_{220}+8z_{8} z_{40}+24z_{8} z_{40} z_{220}+8z_{8} z_{44}\\
&\quad{}+24z_{8} z_{44} z_{220}+16z_{8} z_{65}+48z_{8} z_{65} z_{220}+8z_{8} z_{66}+24z_{8} z_{66} z_{220}\\
&\quad{}-12z_{8} z_{220}+4z_{9}^{2}+12z_{9}^{2} z_{220}+16z_{9} z_{19}+48z_{9} z_{19} z_{220}+8z_{9} z_{22}\\
&\quad{}+24z_{9} z_{22} z_{220}+16z_{9} z_{24}+48z_{9} z_{24} z_{220}+8z_{9} z_{27}+24z_{9} z_{27} z_{220}\\
&\quad{}+8z_{9} z_{29}+24z_{9} z_{29} z_{220}+16z_{9} z_{30}+48z_{9} z_{30} z_{220}+8z_{9} z_{32}\\
&\quad{}+24z_{9} z_{32} z_{220}+8z_{9} z_{34}+24z_{9} z_{34} z_{220}+16z_{9} z_{36}+48z_{9} z_{36} z_{220}\\
&\quad{}+8z_{9} z_{37}+24z_{9} z_{37} z_{220}+8z_{9} z_{40}+24z_{9} z_{40} z_{220}+8z_{9} z_{44}\\
&\quad{}+24z_{9} z_{44} z_{220}+16z_{9} z_{65}+48z_{9} z_{65} z_{220}+8z_{9} z_{66}+24z_{9} z_{66} z_{220}\\
&\quad{}+3z_{10}^{2}+6z_{10}^{2} z_{220}+3z_{11}^{2}+6z_{11}^{2} z_{220}+3z_{12}^{2}+6z_{12}^{2} z_{220}+3z_{13}^{2}\\
&\quad{}+6z_{13}^{2} z_{220}+3z_{14}^{2}+6z_{14}^{2} z_{220}+3z_{15}^{2}+6z_{15}^{2} z_{220}+3z_{16}^{2}\\
&\quad{}+6z_{16}^{2} z_{220}+3z_{17}^{2}+6z_{17}^{2} z_{220}+z_{19}+32z_{19}^{2}+96z_{19}^{2} z_{220}\\
&\quad{}+32z_{19} z_{22}+96z_{19} z_{22} z_{220}+64z_{19} z_{24}+192z_{19} z_{24} z_{220}+32z_{19} z_{27}\\
&\quad{}+96z_{19} z_{27} z_{220}+32z_{19} z_{29}+96z_{19} z_{29} z_{220}+64z_{19} z_{30}+192z_{19} z_{30} z_{220}\\
&\quad{}+32z_{19} z_{32}+96z_{19} z_{32} z_{220}+32z_{19} z_{34}+96z_{19} z_{34} z_{220}+64z_{19} z_{36}\\
&\quad{}+192z_{19} z_{36} z_{220}+32z_{19} z_{37}+96z_{19} z_{37} z_{220}+32z_{19} z_{40}+96z_{19} z_{40} z_{220}\\
&\quad{}+32z_{19} z_{44}+96z_{19} z_{44} z_{220}+64z_{19} z_{65}+192z_{19} z_{65} z_{220}+32z_{19} z_{66}\\
&\quad{}+96z_{19} z_{66} z_{220}+8z_{22}^{2}+24z_{22}^{2} z_{220}+32z_{22} z_{24}+96z_{22} z_{24} z_{220}\\
&\quad{}+16z_{22} z_{27}+48z_{22} z_{27} z_{220}+16z_{22} z_{29}+48z_{22} z_{29} z_{220}+32z_{22} z_{30}\\
&\quad{}+96z_{22} z_{30} z_{220}+16z_{22} z_{32}+48z_{22} z_{32} z_{220}+16z_{22} z_{34}+48z_{22} z_{34} z_{220}\\
&\quad{}+32z_{22} z_{36}+96z_{22} z_{36} z_{220}+16z_{22} z_{37}+48z_{22} z_{37} z_{220}+16z_{22} z_{40}\\
&\quad{}+48z_{22} z_{40} z_{220}+16z_{22} z_{44}+48z_{22} z_{44} z_{220}+32z_{22} z_{65}+96z_{22} z_{65} z_{220}\\
&\quad{}+16z_{22} z_{66}+48z_{22} z_{66} z_{220}+32z_{24}^{2}+96z_{24}^{2} z_{220}+32z_{24} z_{27}\\
&\quad{}+96z_{24} z_{27} z_{220}+32z_{24} z_{29}+96z_{24} z_{29} z_{220}+64z_{24} z_{30}+192z_{24} z_{30} z_{220}\\
&\quad{}+32z_{24} z_{32}+96z_{24} z_{32} z_{220}+32z_{24} z_{34}+96z_{24} z_{34} z_{220}+64z_{24} z_{36}\\
&\quad{}+192z_{24} z_{36} z_{220}+32z_{24} z_{37}+96z_{24} z_{37} z_{220}+32z_{24} z_{40}+96z_{24} z_{40} z_{220}\\
&\quad{}+32z_{24} z_{44}+96z_{24} z_{44} z_{220}+64z_{24} z_{65}+192z_{24} z_{65} z_{220}+32z_{24} z_{66}\\
&\quad{}+96z_{24} z_{66} z_{220}+8z_{27}^{2}+24z_{27}^{2} z_{220}+16z_{27} z_{29}+48z_{27} z_{29} z_{220}\\
&\quad{}+32z_{27} z_{30}+96z_{27} z_{30} z_{220}+16z_{27} z_{32}+48z_{27} z_{32} z_{220}+16z_{27} z_{34}\\
&\quad{}+48z_{27} z_{34} z_{220}+32z_{27} z_{36}+96z_{27} z_{36} z_{220}+16z_{27} z_{37}+48z_{27} z_{37} z_{220}\\
&\quad{}+16z_{27} z_{40}+48z_{27} z_{40} z_{220}+16z_{27} z_{44}+48z_{27} z_{44} z_{220}+32z_{27} z_{65}\\
&\quad{}+96z_{27} z_{65} z_{220}+16z_{27} z_{66}+48z_{27} z_{66} z_{220}+8z_{29}^{2}+24z_{29}^{2} z_{220}\\
&\quad{}+32z_{29} z_{30}+96z_{29} z_{30} z_{220}+16z_{29} z_{32}+48z_{29} z_{32} z_{220}+16z_{29} z_{34}\\
&\quad{}+48z_{29} z_{34} z_{220}+32z_{29} z_{36}+96z_{29} z_{36} z_{220}+16z_{29} z_{37}+48z_{29} z_{37} z_{220}\\
&\quad{}+16z_{29} z_{40}+48z_{29} z_{40} z_{220}+16z_{29} z_{44}+48z_{29} z_{44} z_{220}+32z_{29} z_{65}\\
&\quad{}+96z_{29} z_{65} z_{220}+16z_{29} z_{66}+48z_{29} z_{66} z_{220}+32z_{30}^{2}+96z_{30}^{2} z_{220}\\
&\quad{}+32z_{30} z_{32}+96z_{30} z_{32} z_{220}+32z_{30} z_{34}+96z_{30} z_{34} z_{220}+64z_{30} z_{36}\\
&\quad{}+192z_{30} z_{36} z_{220}+32z_{30} z_{37}+96z_{30} z_{37} z_{220}+32z_{30} z_{40}+96z_{30} z_{40} z_{220}\\
&\quad{}+32z_{30} z_{44}+96z_{30} z_{44} z_{220}+64z_{30} z_{65}+192z_{30} z_{65} z_{220}+32z_{30} z_{66}\\
&\quad{}+96z_{30} z_{66} z_{220}+8z_{32}^{2}+24z_{32}^{2} z_{220}+16z_{32} z_{34}+48z_{32} z_{34} z_{220}\\
&\quad{}+32z_{32} z_{36}+96z_{32} z_{36} z_{220}+16z_{32} z_{37}+48z_{32} z_{37} z_{220}+16z_{32} z_{40}\\
&\quad{}+48z_{32} z_{40} z_{220}+16z_{32} z_{44}+48z_{32} z_{44} z_{220}+32z_{32} z_{65}+96z_{32} z_{65} z_{220}\\
&\quad{}+16z_{32} z_{66}+48z_{32} z_{66} z_{220}+8z_{34}^{2}+24z_{34}^{2} z_{220}+32z_{34} z_{36}\\
&\quad{}+96z_{34} z_{36} z_{220}+16z_{34} z_{37}+48z_{34} z_{37} z_{220}+16z_{34} z_{40}+48z_{34} z_{40} z_{220}\\
&\quad{}+16z_{34} z_{44}+48z_{34} z_{44} z_{220}+32z_{34} z_{65}+96z_{34} z_{65} z_{220}+16z_{34} z_{66}\\
&\quad{}+48z_{34} z_{66} z_{220}+32z_{36}^{2}+96z_{36}^{2} z_{220}+32z_{36} z_{37}+96z_{36} z_{37} z_{220}\\
&\quad{}+32z_{36} z_{40}+96z_{36} z_{40} z_{220}+32z_{36} z_{44}+96z_{36} z_{44} z_{220}+64z_{36} z_{65}\\
&\quad{}+192z_{36} z_{65} z_{220}+32z_{36} z_{66}+96z_{36} z_{66} z_{220}+8z_{37}^{2}+24z_{37}^{2} z_{220}\\
&\quad{}+16z_{37} z_{40}+48z_{37} z_{40} z_{220}+16z_{37} z_{44}+48z_{37} z_{44} z_{220}+32z_{37} z_{65}\\
&\quad{}+96z_{37} z_{65} z_{220}+16z_{37} z_{66}+48z_{37} z_{66} z_{220}+8z_{40}^{2}+24z_{40}^{2} z_{220}\\
&\quad{}+16z_{40} z_{44}+48z_{40} z_{44} z_{220}+32z_{40} z_{65}+96z_{40} z_{65} z_{220}+16z_{40} z_{66}\\
&\quad{}+48z_{40} z_{66} z_{220}+8z_{44}^{2}+24z_{44}^{2} z_{220}+32z_{44} z_{65}+96z_{44} z_{65} z_{220}\\
&\quad{}+16z_{44} z_{66}+48z_{44} z_{66} z_{220}+32z_{65}^{2}+96z_{65}^{2} z_{220}+32z_{65} z_{66}\\
&\quad{}+96z_{65} z_{66} z_{220}+8z_{66}^{2}+24z_{66}^{2} z_{220}+3z_{103}^{2}+6z_{103}^{2} z_{220}+3z_{121}^{2}\\
&\quad{}+6z_{121}^{2} z_{220}+3z_{127}^{2}+6z_{127}^{2} z_{220}-z_{145}^{2}-z_{153}^{2}+2z_{158}-z_{158}^{2}\\
&\quad{}-z_{161}^{2}-z_{164}^{2}-z_{167}^{2}-z_{170}^{2}-z_{173}^{2}-z_{177}^{2}-z_{180}^{2}-z_{183}^{2}-z_{187}^{2}\\
&\quad{}-z_{190}^{2}+z_{193}^{2}+3z_{195}^{2}+6z_{195}^{2} z_{220}-z_{214}^{2}-z_{215}^{2}-z_{216}-z_{216}^{2}\\
&\quad{}+186z_{220}+324z_{220}^{2}+216z_{220}^{3}+24z_{220} z_{225}+24z_{220} z_{225}^{2}+12z_{225}+12z_{225}^{2}\\[3pt]
z_{142}&=-z_{139}+2z_{139} z_{141}\\[3pt]
z_{143}&=z_{139} z_{142}+z_{142}\\[3pt]
z_{154}&=z_{140} z_{153}+7z_{153}\\[3pt]
z_{159}&=-7-z_{140}+z_{140} z_{158}+7z_{158}\\[3pt]
z_{160}&=-1+z_{111}+z_{140} z_{161}+3z_{161}-6z_{161} z_{220}\\[3pt]
z_{162}&=z_{140} z_{161}+7z_{161}\\[3pt]
z_{165}&=z_{140} z_{164}+7z_{164}
\end{align*}
\begin{align*}
z_{168}&=z_{140} z_{167}+7z_{167}\\[3pt]
z_{171}&=z_{140} z_{170}+7z_{170}\\[3pt]
z_{174}&=z_{140} z_{173}+7z_{173}\\[3pt]
z_{172}&=-1+z_{121}+z_{155} z_{175}-3z_{173}-6z_{173} z_{220}+z_{174}\\[3pt]
z_{176}&=-1+z_{124}+z_{140} z_{177}+3z_{177}-6z_{177} z_{220}\\[3pt]
z_{178}&=z_{140} z_{177}+7z_{177}\\[3pt]
z_{181}&=z_{140} z_{180}+7z_{180}\\[3pt]
z_{186}&=-1+z_{130} z_{131}+z_{132}+z_{140} z_{187}+3z_{187}-6z_{187} z_{220}
\end{align*}
\begin{align*}
z_{188}&=z_{140} z_{187}+7z_{187}\\[3pt]
z_{191}&=z_{140} z_{190}+7z_{190}\\[3pt]
z_{192}&=-1+z_{136}+z_{140} z_{193}+3z_{193}-6z_{193} z_{220}\\[3pt]
z_{194}&=z_{140} z_{193}+7z_{193}\\[3pt]
z_{199}&=1+z_{192}+z_{197} z_{205}\\[3pt]
z_{200}&=z_{8}+z_{9}+4z_{19}+2z_{22}+4z_{24}+2z_{27}+2z_{29}+4z_{30}+2z_{32}+2z_{34}+4z_{36}+2z_{37}+2z_{40}+2z_{44}\\*
&\quad{}+4z_{65}+2z_{66}+4z_{193} z_{206}\\[3pt]
z_{207}&=z_{140} z_{198}+7z_{198}\\[3pt]
z_{208}&=z_{200}+z_{200} z_{201}
\end{align*}
\subsection*{Retained residuals}
\begin{align*}
R_{0}&=-1-z_{2} z_{5}+z_{6} z_{7}-2z_{225}
\end{align*}
\begin{align*}
R_{1}&=1+z_{9}-z_{9}^{2}-4z_{9} z_{10}-4z_{9} z_{47}-4z_{10}^{2}-8z_{10} z_{47}-4z_{47}^{2}+z_{48}
\end{align*}
\begin{align*}
R_{2}&=-9+7z_{8}-z_{8}^{2}-8z_{8} z_{45}-4z_{8} z_{52}+24z_{45}-16z_{45}^{2}-16z_{45} z_{52}+12z_{52}-4z_{52}^{2}\\*
&\quad{}+z_{53}
\end{align*}
\begin{align*}
R_{3}&=z_{10}-z_{10}^{2}-8z_{10} z_{59}-4z_{10} z_{61}-16z_{59}^{2}-16z_{59} z_{61}-4z_{61}^{2}+z_{62}
\end{align*}
\begin{align*}
R_{4}&=z_{12}-z_{13}-2z_{34} z_{67}+6z_{67}-z_{72}
\end{align*}
\begin{align*}
R_{5}&=-2z_{1}^{2}+2z_{1} z_{12}+2z_{1} z_{13}+2z_{1} z_{99}+2z_{1} z_{100}-z_{12}^{2}-2z_{12} z_{99}-z_{13}^{2}\\*
&\quad{}-2z_{13} z_{100}+z_{17} z_{101}-z_{78} z_{97}-z_{80} z_{97}-z_{99}^{2}-z_{100}^{2}
\end{align*}
\begin{align*}
R_{6}&=-1+z_{142}^{2}-z_{143}^{2}+z_{223}^{2}
\end{align*}
\begin{align*}
R_{7}&=2z_{144}+z_{144}^{2}+z_{145}^{2}-z_{146}^{2}
\end{align*}
\begin{align*}
R_{8}&=2z_{152}+z_{152}^{2}+z_{153}^{2}-z_{154}^{2}
\end{align*}
\begin{align*}
R_{9}&=-39+z_{130} z_{226}-8z_{140}-12z_{140} z_{220}-36z_{220}+36z_{220}^{2}
\end{align*}
\begin{align*}
R_{10}&=1+2z_{157}+z_{157}^{2}-2z_{158}+z_{158}^{2}-z_{159}^{2}
\end{align*}
\begin{align*}
R_{11}&=2z_{160}+z_{160}^{2}+z_{161}^{2}-z_{162}^{2}
\end{align*}
\begin{align*}
R_{12}&=2z_{163}+z_{163}^{2}+z_{164}^{2}-z_{165}^{2}
\end{align*}
\begin{align*}
R_{13}&=2z_{166}+z_{166}^{2}+z_{167}^{2}-z_{168}^{2}
\end{align*}
\begin{align*}
R_{14}&=2z_{169}+z_{169}^{2}+z_{170}^{2}-z_{171}^{2}
\end{align*}
\begin{align*}
R_{15}&=2z_{172}+z_{172}^{2}+z_{173}^{2}-z_{174}^{2}
\end{align*}
\begin{align*}
R_{16}&=2z_{176}+z_{176}^{2}+z_{177}^{2}-z_{178}^{2}
\end{align*}
\begin{align*}
R_{17}&=2z_{179}+z_{179}^{2}+z_{180}^{2}-z_{181}^{2}
\end{align*}
\begin{align*}
R_{18}&=-1+2z_{182}+z_{182}^{2}+z_{183}^{2}-2z_{184}-z_{184}^{2}
\end{align*}
\begin{align*}
R_{19}&=z_{127}+z_{155} z_{185}-z_{182}-3z_{183}-6z_{183} z_{220}+z_{184}
\end{align*}
\begin{align*}
R_{20}&=2z_{186}+z_{186}^{2}+z_{187}^{2}-z_{188}^{2}
\end{align*}
\begin{align*}
R_{21}&=2z_{189}+z_{189}^{2}+z_{190}^{2}-z_{191}^{2}
\end{align*}
\begin{align*}
R_{22}&=2z_{192}+z_{192}^{2}+z_{193}^{2}-z_{194}^{2}
\end{align*}
\begin{align*}
R_{23}&=-1+z_{197}^{2}+z_{198}^{2}-z_{207}^{2}
\end{align*}
\begin{align*}
R_{24}&=-1+z_{199}^{2}+z_{200}^{2}-z_{208}^{2}
\end{align*}
\begin{align*}
R_{25}&=-2+2z_{0}^{2}+2z_{0} z_{4}-2z_{0} z_{145}-2z_{0} z_{153}+z_{1}^{2}-2z_{1} z_{183}+z_{2}^{2}-2z_{2} z_{173}\\*
&\quad{}+z_{4}^{2}-2z_{4} z_{153}+z_{8}^{2}-2z_{8} z_{158}-2z_{8} z_{193}+z_{9}^{2}-2z_{9} z_{161}-2z_{9} z_{193}\\*
&\quad{}+z_{10}^{2}-2z_{10} z_{164}+z_{11}^{2}-2z_{11} z_{167}+z_{12}^{2}-2z_{12} z_{170}+z_{13}^{2}-2z_{13} z_{177}\\*
&\quad{}+z_{14}^{2}-2z_{14} z_{180}+z_{15}^{2}-2z_{15} z_{187}+z_{16}^{2}-2z_{16} z_{190}-8z_{19} z_{193}\\*
&\quad{}-4z_{22} z_{193}-8z_{24} z_{193}-4z_{27} z_{193}-4z_{29} z_{193}-8z_{30} z_{193}-4z_{32} z_{193}\\*
&\quad{}-4z_{34} z_{193}-8z_{36} z_{193}-4z_{37} z_{193}-4z_{40} z_{193}-4z_{44} z_{193}-8z_{65} z_{193}\\*
&\quad{}-4z_{66} z_{193}-z_{140} z_{210}+z_{144}^{2}+z_{152}^{2}+z_{157}^{2}+2z_{158}+z_{160}^{2}+z_{163}^{2}\\*
&\quad{}+z_{166}^{2}+z_{169}^{2}+z_{172}^{2}+z_{176}^{2}+z_{179}^{2}+z_{182}^{2}+z_{186}^{2}+z_{189}^{2}+z_{192}^{2}\\*
&\quad{}+2z_{193}^{2}-6z_{210}-z_{214}^{2}-z_{215}^{2}-z_{216}-z_{216}^{2}
\end{align*}
\begin{align*}
R_{26}&=-z_{103} z_{211}+z_{127}
\end{align*}
\begin{align*}
R_{27}&=-z_{108} z_{212}+z_{127}
\end{align*}
\begin{align*}
R_{28}&=-1-8z_{8} z_{19}-4z_{8} z_{22}-8z_{8} z_{24}-4z_{8} z_{27}-4z_{8} z_{29}-8z_{8} z_{30}-4z_{8} z_{32}\\*
&\quad{}-4z_{8} z_{34}-8z_{8} z_{36}-4z_{8} z_{37}-4z_{8} z_{40}-4z_{8} z_{44}-8z_{8} z_{65}-4z_{8} z_{66}\\*
&\quad{}-8z_{9} z_{19}-4z_{9} z_{22}-8z_{9} z_{24}-4z_{9} z_{27}-4z_{9} z_{29}-8z_{9} z_{30}-4z_{9} z_{32}\\*
&\quad{}-4z_{9} z_{34}-8z_{9} z_{36}-4z_{9} z_{37}-4z_{9} z_{40}-4z_{9} z_{44}-8z_{9} z_{65}-4z_{9} z_{66}-z_{19}\\*
&\quad{}-16z_{19}^{2}-16z_{19} z_{22}-32z_{19} z_{24}-16z_{19} z_{27}-16z_{19} z_{29}-32z_{19} z_{30}\\*
&\quad{}-16z_{19} z_{32}-16z_{19} z_{34}-32z_{19} z_{36}-16z_{19} z_{37}-16z_{19} z_{40}-16z_{19} z_{44}\\*
&\quad{}-32z_{19} z_{65}-16z_{19} z_{66}-4z_{22}^{2}-16z_{22} z_{24}-8z_{22} z_{27}-8z_{22} z_{29}-16z_{22} z_{30}\\*
&\quad{}-8z_{22} z_{32}-8z_{22} z_{34}-16z_{22} z_{36}-8z_{22} z_{37}-8z_{22} z_{40}-8z_{22} z_{44}-16z_{22} z_{65}\\*
&\quad{}-8z_{22} z_{66}-16z_{24}^{2}-16z_{24} z_{27}-16z_{24} z_{29}-32z_{24} z_{30}-16z_{24} z_{32}-16z_{24} z_{34}\\*
&\quad{}-32z_{24} z_{36}-16z_{24} z_{37}-16z_{24} z_{40}-16z_{24} z_{44}-32z_{24} z_{65}-16z_{24} z_{66}-4z_{27}^{2}\\*
&\quad{}-8z_{27} z_{29}-16z_{27} z_{30}-8z_{27} z_{32}-8z_{27} z_{34}-16z_{27} z_{36}-8z_{27} z_{37}-8z_{27} z_{40}\\*
&\quad{}-8z_{27} z_{44}-16z_{27} z_{65}-8z_{27} z_{66}-4z_{29}^{2}-16z_{29} z_{30}-8z_{29} z_{32}-8z_{29} z_{34}\\*
&\quad{}-16z_{29} z_{36}-8z_{29} z_{37}-8z_{29} z_{40}-8z_{29} z_{44}-16z_{29} z_{65}-8z_{29} z_{66}-16z_{30}^{2}\\*
&\quad{}-16z_{30} z_{32}-16z_{30} z_{34}-32z_{30} z_{36}-16z_{30} z_{37}-16z_{30} z_{40}-16z_{30} z_{44}\\*
&\quad{}-32z_{30} z_{65}-16z_{30} z_{66}-4z_{32}^{2}-8z_{32} z_{34}-16z_{32} z_{36}-8z_{32} z_{37}-8z_{32} z_{40}\\*
&\quad{}-8z_{32} z_{44}-16z_{32} z_{65}-8z_{32} z_{66}-4z_{34}^{2}-16z_{34} z_{36}-8z_{34} z_{37}-8z_{34} z_{40}\\*
&\quad{}-8z_{34} z_{44}-16z_{34} z_{65}-8z_{34} z_{66}-16z_{36}^{2}-16z_{36} z_{37}-16z_{36} z_{40}-16z_{36} z_{44}\\*
&\quad{}-32z_{36} z_{65}-16z_{36} z_{66}-4z_{37}^{2}-8z_{37} z_{40}-8z_{37} z_{44}-16z_{37} z_{65}-8z_{37} z_{66}\\*
&\quad{}-4z_{40}^{2}-8z_{40} z_{44}-16z_{40} z_{65}-8z_{40} z_{66}-4z_{44}^{2}-16z_{44} z_{65}-8z_{44} z_{66}\\*
&\quad{}-16z_{65}^{2}-16z_{65} z_{66}-4z_{66}^{2}-z_{145}^{2}-z_{153}^{2}+2z_{158}-z_{158}^{2}-z_{161}^{2}\\*
&\quad{}-z_{164}^{2}-z_{167}^{2}-z_{170}^{2}-z_{173}^{2}-z_{177}^{2}-z_{180}^{2}-z_{183}^{2}-z_{187}^{2}-z_{190}^{2}\\*
&\quad{}+z_{193}^{2}-z_{214}^{2}-z_{215}^{2}-z_{216}-z_{216}^{2}
\end{align*}
\begin{align*}
R_{29}&=-1-4z_{105}^{2}-4z_{110}^{2}-4z_{113}^{2}-4z_{116}^{2}-4z_{119}^{2}-4z_{123}^{2}-4z_{126}^{2}+z_{127}^{2}\\*
&\quad{}-4z_{129}^{2}-4z_{132}^{2}-4z_{135}^{2}-4z_{138}^{2}-z_{217}^{2}-z_{218}^{2}-z_{219}-z_{219}^{2}
\end{align*}
\begin{align*}
R_{30}&=5+z_{140}+8z_{141}-8z_{141}^{2}-z_{223}^{2}
\end{align*}
\endgroup
The required equation is therefore exactly
\[D_{83}(x_0,\ldots,x_{82})=\sum_{j=0}^{30}R_j^2=0.\]
The electronic source preserves the source-row indices, unit-pivot witnesses and binary graph, so these printed definitions can be checked independently against the literal coefficient files.

% END retained input: P83_definition.tex


\clearpage
\begin{thebibliography}{9}
\bibitem{source172}
ProveIt, package 172, \emph{Pooled Powers, Round 21}, positive-ten component,
25 October 2026. Frozen proof and equation retained in this archive.
\url{https://github.com/VladimirReshetnikov/ProveIt/tree/cdfea7d90176775d74c1c30afbe49fa78ade9897/Computability/DiophantineConsistency/con-zfc-diophantine-forms/172-pooled-powers-round21}.
\bibitem{source178}
ProveIt, package 178, \emph{Polynomial Centers for Diophantine Consistency},
25 October 2026. Whole-product center and fixed coefficient recipe.
\url{https://github.com/VladimirReshetnikov/ProveIt/tree/cdfea7d90176775d74c1c30afbe49fa78ade9897/Computability/DiophantineConsistency/con-zfc-diophantine-forms/178-polynomial-centers}.
\bibitem{incomingAffine}
ProveIt incoming archive, \emph{Affine Horizons and Centered Certificates},
file \path{ConZFC_Affine_Horizons_Centered_Certificates_2026-10-26.zip},
at the pinned snapshot. The filename's date is retained as source metadata.
\url{https://github.com/VladimirReshetnikov/ProveIt/tree/cdfea7d90176775d74c1c30afbe49fa78ade9897/docs/incoming}.
\bibitem{source166}
ProveIt, package 166, \emph{Fixed-Base Modular Construction}, including
common modular numerators, frame, and C222 calibration sources.
\url{https://github.com/VladimirReshetnikov/ProveIt/tree/cdfea7d90176775d74c1c30afbe49fa78ade9897/Computability/DiophantineConsistency/con-zfc-diophantine-forms/166-fixed-base-modular}.
\bibitem{incomingCofinal}
ProveIt incoming archive, \emph{Cofinal Pell Compression},
\path{ConZFC_Cofinal_Pell_Compression_2026-10-25.zip},
at the same snapshot. One scaled square and exact source replay.
\bibitem{source177}
ProveIt, package 177, \emph{Oriented Power Pooling and Polynomial Root
Centering}, 25 October 2026. Direct profiles and acyclic unit elimination.
\url{https://github.com/VladimirReshetnikov/ProveIt/tree/cdfea7d90176775d74c1c30afbe49fa78ade9897/Computability/DiophantineConsistency/con-zfc-diophantine-forms/177-oriented-powers}.
\bibitem{source175}
ProveIt, package 175, \emph{Exact Integer Projections for Diophantine
Consistency}, 25 October 2026. The \(92\)-integer, degree-\(692\) comparison.
\url{https://github.com/VladimirReshetnikov/ProveIt/tree/cdfea7d90176775d74c1c30afbe49fa78ade9897/Computability/DiophantineConsistency/con-zfc-diophantine-forms/175-exact-contractions}.
\end{thebibliography}
\end{document}
